\documentclass{article}
\usepackage{fancyhdr}
\usepackage[american]{babel}

\usepackage{amssymb, amsmath, amsthm, amsfonts, enumerate}
\usepackage{amsmath,setspace,scalefnt}
\usepackage[dvipsnames,svgnames,table]{xcolor}
\usepackage{graphicx,tikz,caption,subcaption}
\usetikzlibrary{patterns}
\usetikzlibrary{shapes}
\usepackage{fullpage}
\usepackage{hyperref}
%\usepackage[notref,notcite]{showkeys}
\usepackage{enumitem,verbatim}%,itemize
\usepackage{array,booktabs,longtable}
\usepackage{multicol}
\usepackage{float}
\usepackage{aliascnt}
\usepackage{cleveref}
%\usepackage{ulem}
\usepackage[normalem]{ulem}
%\usepackage{lineno}
\usepackage{tikz-cd}
\usepackage{xparse}
\usepackage{authblk}
\renewcommand{\Affilfont}{\small}


\usepackage[margin=2.5cm]{geometry}

% Keep the many multiline formulas visually open while avoiding excess vertical space.
\AtBeginDocument{%
  \addtolength{\abovedisplayskip}{-1pt}%
  \addtolength{\belowdisplayskip}{-1pt}%
  \addtolength{\belowdisplayshortskip}{-1pt}%
}

\definecolor{gray}{rgb}{0.25, 0.25, 0.25}

\newtheorem{theorem}{Theorem}[section]
\newaliascnt{lemma}{theorem}
\newtheorem{lemma}[lemma]{Lemma}
\aliascntresetthe{lemma}
\newaliascnt{cor}{theorem}
\newtheorem{cor}[cor]{Corollary}
\aliascntresetthe{cor}
\newaliascnt{conj}{theorem}
\newtheorem{conj}[conj]{Conjecture}
\aliascntresetthe{conj}
\newaliascnt{question}{theorem}
\newtheorem{question}[question]{Question}
\aliascntresetthe{question}
\newtheorem*{theorem*}{Theorem}

\theoremstyle{definition}
\newtheorem{subclaim}{Subclaim}
\DeclareTextCompositeCommand{\v}{OT1}{l}{l\nobreak\hspace{-.1em}'} %for the name of Dan Kral

\global\long\def\si#1{\textcolor{red}{\textbf{[SI comments:} #1\textbf{]}}}
\global\long\def\sl#1{\textcolor{blue}{\textbf{[SL comments:} #1\textbf{]}}}
\global\long\def\hy#1{\textcolor{purple}{\textbf{[HY comments:} #1\textbf{]}}}
%define your comment if you need. 
%define your comment if you need. 

\newcommand{\eps}{\varepsilon}
\newcommand{\pc}{p_{\mathrm c}}
\newcommand{\W}{\mathcal W}
\newcommand{\what}{\widehat{\mathcal W}}
\newcommand{\dd}{\,d}
\newcommand{\EE}{\mathbb E}
\newcommand{\RR}{\mathbb R}
\newcommand{\1}{\mathbf 1}
\newcommand{\cut}{\square}
\newcommand{\dist}{\operatorname{dist}}
\newcommand{\argmin}{\operatorname*{argmin}}
\newcommand{\argmax}{\operatorname*{argmax}}
\newcommand{\jp}{J_p}
\newcommand{\ip}{I_p}
\newcommand{\sH}{S_H}
\newcommand{\PhiH}{\Phi_H}
\newcommand{\zr}{\zeta_d}
\newcommand{\sm}{s_{\mathrm M}}
\newcommand{\ER}{\mathrm{ER}}


\newcounter{propcounter}
\newenvironment{proplist}[1][]{%
    \stepcounter{propcounter}%
    \enumerate[label = {\bfseries \Alph{propcounter}\arabic{enumi}}, #1]%
}
{\endenumerate}


\newenvironment{poc}{\begin{proof}[Proof of claim]\renewcommand*{\qedsymbol}{$\blacksquare$}}{\end{proof}}
\newenvironment{posc}{\begin{proof}[Proof of subclaim]\renewcommand*{\qedsymbol}{$\blacksquare$}}{\end{proof}}



\theoremstyle{plain}
\newaliascnt{claim}{theorem}
\newtheorem{claim}[claim]{Claim}
\aliascntresetthe{claim}

\newaliascnt{op}{theorem}
\newtheorem{op}[op]{Open Question}
\aliascntresetthe{op}
\newaliascnt{prob}{theorem}
\newtheorem{prob}[prob]{Problem}
\aliascntresetthe{prob}
\newaliascnt{prop}{theorem}
\newtheorem{prop}[prop]{Proposition}
\aliascntresetthe{prop}
\newaliascnt{proposition}{theorem}
\newtheorem{proposition}[proposition]{Proposition}
\aliascntresetthe{proposition}
\newaliascnt{corollary}{theorem}
\newtheorem{corollary}[corollary]{Corollary}
\aliascntresetthe{corollary}
\theoremstyle{definition}
\newaliascnt{rmk}{theorem}
\newtheorem{rmk}[rmk]{Remark}
\aliascntresetthe{rmk}
\newaliascnt{remark}{theorem}
\newtheorem{remark}[remark]{Remark}
\aliascntresetthe{remark}
\theoremstyle{definition}
\newaliascnt{intp}{theorem}
\newtheorem{intp}[intp]{Interpretation}
\aliascntresetthe{intp}
\theoremstyle{definition}
\newaliascnt{kn}{theorem}
\newtheorem{kn}[kn]{Known}
\aliascntresetthe{kn}
\theoremstyle{definition}
\newaliascnt{defn}{theorem}
\newtheorem{defn}[defn]{Definition}
\aliascntresetthe{defn}
\theoremstyle{definition}
\newaliascnt{expl}{theorem}
\newtheorem{expl}[expl]{Example}
\aliascntresetthe{expl}
\theoremstyle{definition}
\newaliascnt{obs}{theorem}
\newtheorem{obs}[obs]{Observation}
\aliascntresetthe{obs}
\crefname{theorem}{Theorem}{Theorems}
\Crefname{theorem}{Theorem}{Theorems}
\crefname{lemma}{Lemma}{Lemmas}
\Crefname{lemma}{Lemma}{Lemmas}
\crefname{cor}{Corollary}{Corollaries}
\Crefname{cor}{Corollary}{Corollaries}
\crefname{corollary}{Corollary}{Corollaries}
\Crefname{corollary}{Corollary}{Corollaries}
\crefname{conj}{Conjecture}{Conjectures}
\Crefname{conj}{Conjecture}{Conjectures}
\crefname{question}{Question}{Questions}
\Crefname{question}{Question}{Questions}
\crefname{claim}{Claim}{Claims}
\Crefname{claim}{Claim}{Claims}
\crefname{op}{Open Question}{Open Questions}
\Crefname{op}{Open Question}{Open Questions}
\crefname{prob}{Problem}{Problems}
\Crefname{prob}{Problem}{Problems}
\crefname{prop}{Proposition}{Propositions}
\Crefname{prop}{Proposition}{Propositions}
\crefname{proposition}{Proposition}{Propositions}
\Crefname{proposition}{Proposition}{Propositions}
\crefname{rmk}{Remark}{Remarks}
\Crefname{rmk}{Remark}{Remarks}
\crefname{remark}{Remark}{Remarks}
\Crefname{remark}{Remark}{Remarks}
\crefname{intp}{Interpretation}{Interpretations}
\Crefname{intp}{Interpretation}{Interpretations}
\crefname{kn}{Known}{Knowns}
\Crefname{kn}{Known}{Knowns}
\crefname{defn}{Definition}{Definitions}
\Crefname{defn}{Definition}{Definitions}
\crefname{expl}{Example}{Examples}
\Crefname{expl}{Example}{Examples}
\crefname{obs}{Observation}{Observations}
\Crefname{obs}{Observation}{Observations}
\crefname{subclaim}{Subclaim}{Subclaims}
\Crefname{subclaim}{Subclaim}{Subclaims}
\newcommand{\<}{\subseteq}
\newcommand{\blue}[1]{\textcolor[rgb]{0.00,0.00,1.00}{#1}}
\newcommand{\red}[1]{\textcolor[rgb]{1.00,0.00,0.00}{#1}}

\NewDocumentEnvironment{shl}{}
  {%
    \begingroup
    \color{blue}%
    \bfseries
    \boldmath
    [SL:\space\ignorespaces
  }
  {%
    \unskip]\endgroup\ignorespacesafterend
  }


% Notation specific to this paper.
\newcommand{\lo}{\ell}
\newcommand{\HS}{\mathrm{HS}}

\title{Multipodal structure of optimizers in the upper-tail variational problem}

\date{}

\author[1,2]{Sangho Lim}
\author[3]{Seonghyuk Im}
\author[1,2]{Taeyoung Kim}
\author[4]{Kyeongsik Nam}
\author[1,2]{Hongseok Yang}
\affil[1]{School of Computing, KAIST, Daejeon, Korea}
\affil[2]{School of Computational Sciences, Korea Institute for Advanced Study (KIAS), Seoul, Korea}
\affil[3]{Center for AI and Natural Sciences, Korea Institute for Advanced Study (KIAS), Seoul, Korea}
\affil[4]{Department of Mathematical Sciences, KAIST, Daejeon, Korea}
\affil[ ]{E-mail addresses: \texttt{lim.sang@kaist.ac.kr},
\texttt{seonghyuk@kias.re.kr},
\texttt{taeyoungkim21@kaist.ac.kr},
\texttt{ksnam@kaist.ac.kr},
\texttt{hongseokyang@kias.re.kr}}



\begin{document}

\maketitle
\begin{abstract}
Let $H$ be a fixed graph and let $t(H,\cdot)$ denote its homomorphism density.
For fixed $0<p<r<1$, the upper-tail event $t(H,G(n,p))\ge r^{|E(H)|}$ in the
dense Erd\H{o}s--R\'enyi random graph is governed by the Chatterjee--Varadhan
variational problem, whose minimizers describe the typical structure of the
conditioned random graph.  These minimizers are known only in special cases.
We prove that, for every graph $H$ and all parameters, every minimizer is
multipodal, that is, it is the graphon of a stochastic block model with
finitely many blocks.  Moreover, unless $H$ is a perfect matching, the blocks
can be ordered so that the edge density between two blocks is nondecreasing in
the index of each block.
We also consider the graphons of maximal entropy with a prescribed edge
density and a prescribed $H$-density above its Erd\H{o}s--R\'enyi value.
Every such maximizer that satisfies a natural constraint qualification is
multipodal with ordered blocks, and the values of the $H$-density at which
some maximizer violates the constraint qualification form a set of Lebesgue
measure zero.
The proof is based on inserting into an optimizer $W$ a small block of
vertices whose edge probabilities to the existing vertices are arbitrary.
This shows that almost every row $W(x,\cdot)$ of the optimizer minimizes one
and the same functional over all measurable functions from $[0,1]$ to $[0,1]$.  A
strict supermodularity property forces these minimizers to be totally
ordered, and an analysis of the second variation shows that there are only
finitely many of them.
\end{abstract}

\section{Introduction}\label{sec:introduction}

The Erd\H{o}s--R\'enyi graph $G(n,p)$ is the random graph on $n$ vertices in
which each edge is present independently with probability $p$.  For a fixed
graph $H$ and fixed $0<p<r<1$, consider the upper-tail event that the
homomorphism density of $H$ in $G(n,p)$ is at least $r^{|E(H)|}$, which is,
to leading order, its typical value in $G(n,r)$.  Two questions
of Chatterjee and Varadhan~\cite{ChatterjeeVaradhanLDP} guide the study of
this event: how small is its probability, and what does $G(n,p)$ typically
look like conditioned on it?  In the dense regime considered here, both
questions are answered by a variational problem on graphons, and the second
question asks for the structure of its minimizers.

These minimizers are known only in special cases.  If the parameters lie in
the replica-symmetric phase, the unique minimizer is the constant graphon, and
the conditioned graph looks like $G(n,r)$.  Lubetzky and
Zhao~\cite{lubetzky2015large} determined this phase exactly when $H$ is
regular.  In the complementary symmetry-breaking phase, the minimizers are
nonconstant, and Lubetzky and Zhao~\cite[Section~7]{lubetzky2015large} asked
whether they are always bipodal, that is, whether the conditioned graph always
splits into two blocks of vertices with homogeneous edge densities within and
between the blocks.  The companion paper~\cite{LimImKimNamYangBipodal}
confirmed this for $d$-regular $H$ near every point of the phase boundary with
target density $r\ne(d-1)/d$, and along a curve approaching the boundary at the
exceptional target density $(d-1)/d$, using the entropy-maximization results of
Kenyon, Radin, Ren, and Sadun~\cite{kenyon2014entropy}.  For nonregular graphs,
the phase boundary is studied in~\cite{LimImKimNamYangBoundary}, where it is
shown, for instance, that stars with many leaves have nonconstant minimizers on
the boundary itself.  Away from the boundary, and for nonregular graphs, even
the following weaker question has remained open: is every minimizer a graphon
with \emph{finitely many} blocks?

A closely related question arises for dense graphs with prescribed edge
density and prescribed density of $H$.  Their typical structure is given by the
graphons of maximal entropy under these two constraints.  Radin, Ren, and
Sadun~\cite{RadinRenSadun2014} and subsequent
works~\cite{KenyonRadinRenSadun2017b,NeemanRadinSadun2023,NeemanRadinSadun2024,RadinSadun2013,RadinSadun2015}
revealed a rich phase diagram when $H$ is a triangle.  The maximizers found
there, numerically and in the regimes where they have been determined
rigorously, are \emph{multipodal}: they have finitely many blocks.  Kenyon,
Radin, Ren, and Sadun~\cite{KenyonRadinRenSadun2017} proved multipodality when
$H$ is a star, and showed in~\cite{kenyon2014entropy} that for every graph $H$
the maximizers are bipodal when the $H$-density is slightly above the
Erd\H{o}s--R\'enyi value, for all but finitely many edge densities.  Beyond
these regimes, multipodality has not been established.

In this paper we show that every minimizer of the upper-tail variational
problem is multipodal, for every graph $H$ and every $0<p<r<1$, and that the
same holds for the maximizers of the entropy with prescribed edge density and
$H$-density, for almost every value of the $H$-density.  In both settings the
blocks are moreover ordered: the blocks can be listed so that the edge density
between two blocks is nondecreasing in the index of each block.

Finite multipodality is not a formal consequence of the fact that only
finitely many densities are prescribed.  Graphons that are uniquely determined
by finitely many subgraph densities can have a very complex structure; see,
for example,~\cite{CooperKralMartins2018,GrzesikKralLovasz2020}.  Our argument
uses both the entropy and the global optimality of the minimizer in an
essential way.

\subsection{Main results}\label{sec:main-results}

A \emph{graphon} is a symmetric measurable function $W:[0,1]^2\to[0,1]$, and
$\W_0$ denotes the space of graphons~\cite{LovaszGraphLimits}.  For a graph $H$
with $m:=|E(H)|\ge1$ edges and a graphon $W$, the edge density, the
homomorphism density of $H$, and the relative entropy with respect to $p$ are
\begin{equation}\label{eq:graphon-quantities}
  e(W)=\int_{[0,1]^2}W,
  \qquad
  t(H,W)=\int_{[0,1]^{V(H)}}\prod_{ij\in E(H)}W(x_i,x_j)\prod_{i\in V(H)}\dd x_i,
  \qquad
  I_p(W)=\int_{[0,1]^2}J_p(W),
\end{equation}
where $J_p(z)=z\log\frac zp+(1-z)\log\frac{1-z}{1-p}$.  For $0<p<r<1$, the
Chatterjee--Varadhan variational problem is
\begin{equation}\label{eq:graphon-variational-problem}
  \Phi_H(p,r)=\inf\bigl\{I_p(W):W\in\W_0,\ t(H,W)\ge r^m\bigr\}.
\end{equation}
Chatterjee and Varadhan~\cite{ChatterjeeVaradhanLDP} proved that the
probability of the event $t(H,G(n,p))\ge r^m$ is
$\exp(-\Phi_H(p,r)\binom n2+o(n^2))$, and that, conditioned on this event,
$G(n,p)$ is close in cut distance to the set of minimizers
of~\eqref{eq:graphon-variational-problem} with probability tending to one.

A graphon $W$ is \emph{multipodal} if there are a finite measurable partition
$[0,1]=A_1\cup\dots\cup A_M$ and numbers $b_{ij}=b_{ji}\in[0,1]$ such that
$W=b_{ij}$ almost everywhere on $A_i\times A_j$.  Its blocks are
\emph{ordered} if they can be labeled so that $b_{ij}\le b_{i'j}$ whenever
$i\le i'$, for every $j$.  Equivalently, the distinct rows $W(x,\cdot)$ of $W$
are totally ordered pointwise.  In graph terms, an ordered multipodal graphon
describes a stochastic block model in which the blocks can be ranked so that
every vertex in a higher-ranked block is at least as likely as every vertex in
a lower-ranked block to be adjacent to any given vertex.  For example, a
bipodal graphon of the form $f(x)f(y)$, where $f$ takes two values, has ordered
blocks; the optimizers constructed in~\cite{LimImKimNamYangBipodal} near the
exceptional target density $(d-1)/d$ are of this form.

\begin{theorem}[Upper tails]\label{thm:upper-tail}
Let $H$ be a graph with at least one edge, and let $0<p<r<1$.  Every minimizer
of~\eqref{eq:graphon-variational-problem} is multipodal.  If $H$ is not a
perfect matching, its blocks are ordered.  If $H$ is a perfect matching, the
unique minimizer is the constant graphon $W\equiv r$.
\end{theorem}

Theorem~\ref{thm:upper-tail} holds throughout the parameter space, including
the symmetry-breaking phase, where the minimizers are not known explicitly.
It does not bound the number of blocks, and it does not assert uniqueness.

We next consider graphons with a prescribed edge density $\eps\in(0,1)$ and a
prescribed $H$-density $\tau$.  Recall the graphon entropy
\[
  s(W)=\int_{[0,1]^2}S_0\bigl(W(x,y)\bigr)\dd x\dd y,
  \qquad
  S_0(z)=-\frac12\bigl(z\log z+(1-z)\log(1-z)\bigr).
\]
Let $\tau_{\max}(\eps):=\max\{t(H,W):e(W)=\eps\}$, and consider
\begin{equation}\label{eq:entropy-maximization}
  \sup\bigl\{s(W):W\in\W_0,\ e(W)=\eps,\ t(H,W)=\tau\bigr\},
  \qquad
  \eps^m<\tau<\tau_{\max}(\eps).
\end{equation}
The maximizers of~\eqref{eq:entropy-maximization} describe the typical
structure of large graphs with edge density $\eps$ and $H$-density
$\tau$~\cite{RadinRenSadun2014,RadinSadun2013}; see
also~\cite{DemboLubetzky2018} for the underlying large-deviation principle.  The two constraints of
this problem can be degenerate at a maximizer, and we need a constraint
qualification (Definition~\ref{def:constraint-qualification}) that allows the
two densities to be corrected independently by bounded perturbations.

\begin{theorem}[Prescribed edge and subgraph densities]\label{thm:fixed-densities}
Let $H$ be a graph with a vertex of degree at least two, let $\eps\in(0,1)$, and
let $\eps^m<\tau<\tau_{\max}(\eps)$.  Every maximizer
of~\eqref{eq:entropy-maximization} that satisfies the constraint qualification
is multipodal with ordered blocks.
\end{theorem}

\begin{cor}\label{cor:almost-every-tau}
Let $H$ and $\eps$ be as in Theorem~\ref{thm:fixed-densities}.  The set of
$\tau\in(\eps^m,\tau_{\max}(\eps))$ for which~\eqref{eq:entropy-maximization} has
a maximizer violating the constraint qualification has Lebesgue measure zero.
Consequently, for almost every such $\tau$, every maximizer is multipodal with
ordered blocks.
\end{cor}

\subsection{Proof strategy}\label{sec:proof-strategy}

The natural starting point is the Euler--Lagrange equation.  For a minimizer
$W$ of~\eqref{eq:graphon-variational-problem}, it states that
$J_p'(W(x,y))=\mu\,\mathcal D_W(x,y)$ for a Lagrange multiplier $\mu>0$, where
$\mathcal D_W$ is the first variation of $t(H,\cdot)$.  One can rewrite this as
the statement that every row $W_x=W(x,\cdot)$ is a critical point of a
functional $\mathcal F_W$ defined on all possible rows.  However, the set of
critical points of a nonlinear functional on an infinite-dimensional space can
be complicated, and the equation alone does not seem to force finitely many
row types for general $H$.  Our proof instead proceeds in three steps.

\begin{enumerate}[label=(\roman*),leftmargin=2.2em]
\item \emph{Every row is a global minimizer.}  Given an arbitrary candidate row
$u:[0,1]\to[0,1]$, we insert into $W$ a new block of vertices of small measure
$\theta$ whose row is $u$, and we then shift the remaining values slightly so
that the constraint holds exactly.  Comparing the cost of this graphon with that
of $W$ to first order in $\theta$ shows that $\mathcal F_W(u)\ge c_W$ for every $u$,
where $c_W$ is an explicit constant.  The functional $\mathcal F_W$ involves the
\emph{rooted densities} of $H$, which record the density of copies of $H$ that
use the new vertex.  Averaging over the rows of $W$ shows that almost every row
attains the bound $c_W$.  Hence all rows of $W$ are global minimizers of one and
the same functional.  See Lemma~\ref{lem:upper-tail-insertion}.

\item \emph{Minimizers are ordered.}  When two rows $u_1,u_2$ are replaced by
their pointwise minimum and maximum, the sum of the entropy and linear parts of
$\mathcal F_W$ over the two rows does not change, while the sum of the rooted
densities strictly increases whenever $u_1$ and $u_2$ cross; this uses a vertex
of $H$ of degree at least two and the fact that $W$ is bounded below by a
positive constant.
Since the multiplier $\mu$ is positive, two crossing minimizers would produce a
pair of rows with strictly smaller total value, which is impossible.  Hence the
set of global minimizers is totally ordered.  See
Proposition~\ref{prop:row-chain}.

\item \emph{There are finitely many minimizers.}  A totally ordered set of
rows could still be a continuum.  The Euler--Lagrange equation for rows is a
compact fixed-point equation, so the set of global minimizers is compact.  The
same strict supermodularity shows that the second variation of $\mathcal F_W$
is coercive on mean-zero perturbations.  This allows us to parametrize the
critical rows near a minimizer analytically by their mean.  If minimizers
accumulated, the Lagrange multiplier for the mean would vanish on an interval,
and a whole interval of means would consist of minimizers; a compactness
argument rules this out.  Finitely many row types then yield finitely many
blocks.  See Theorem~\ref{thm:row-criterion}.
\end{enumerate}

For the problem~\eqref{eq:entropy-maximization} with two constraints, the
same argument applies once both densities can be restored after inserting a
block, which is what the constraint qualification guarantees.  To prove
Corollary~\ref{cor:almost-every-tau}, we show that a maximizer violating the
constraint qualification forces the left difference quotients of the value
function of the problem to be unbounded, which can only happen on a null set
because the value function is monotone.

\subsection{Related work}

The Euler--Lagrange equations of constrained graphon problems have been used
to establish multipodality when they separate.  For stars, the first variation
of the star density is a sum of a function of $x$ and a function of $y$, and
the equation reduces to a scalar equation for the degree function; this is the
mechanism behind the multipodality results of Kenyon, Radin, Ren, and
Sadun~\cite{KenyonRadinRenSadun2017} for edge and star constraints, where a
uniform bound on the number of blocks is also obtained.  Radin and
Sadun~\cite{RadinSadunEdge2Star} studied the optimal graphons of the
edge--$2$-star model in detail.  For the edge--triangle model, Koch~\cite{Koch2016}
studied an ordering of the vertices of optimal graphons, and Neeman, Radin,
and Sadun~\cite{NeemanRadinSadun2023,NeemanRadinSadun2024} established
detailed properties of several phases.  Radin and
Sadun~\cite{RadinSadunNearExtreme} analyzed optimal graphons for constraints
close to the boundary of the feasible region.
For exponential random graph models with nonnegative parameters, which are
Lagrangian relaxations of problems of the above type, Chatterjee and
Diaconis~\cite{ChatterjeeDiaconis2013} showed that the optimizers are constant;
in the symmetry-breaking phase, however, the minimizers
of~\eqref{eq:graphon-variational-problem} are nonconstant and are not
minimizers of the corresponding Lagrangian.

Alvarado, Ramon, and Wang~\cite{AlvaradoRamonWang2023,AlvaradoRamonWang2026}
consider entropy maximization for multi-relational hypergraphons under
constraints on subgraph densities, and state that for almost every value of
the constraints the optimizers are step functions whose number of steps is
bounded in terms of the number of constraints.  Several steps of their
arguments, however, do not hold.  Representatively, the arguments compare step
functions with $k$ and $k+1$ steps and assert that passing to $k+1$ steps
produces no smaller value of the objective.  This fails in general.  In
Appendix~\ref{app:arw} we give constraints on the densities of the edge, the
$2$-star, and the $3$-star, together with a non-convex objective, for which the
bound on the number of steps fails on a set of constraint values of positive
measure; a second example shows the same for the entropy under constraints of
the more general form allowed in~\cite{AlvaradoRamonWang2026}.  The later
preprint also deduces local optimality in $L^1$ from second-order conditions, although
inserting a small block of vertices is a perturbation that is small in $L^1$
and changes the objective at first order.  These examples do not decide whether
the conclusions of~\cite{AlvaradoRamonWang2023,AlvaradoRamonWang2026} hold for
the entropy under subgraph density constraints.  Our proofs are independent of
these works and follow a different route: we never approximate optimizers by
step functions of bounded size, and the insertion of small blocks is precisely
what gives the first-order condition satisfied by every optimizer.
Appendix~\ref{app:arw} contains the counterexamples and explains why they do
not affect the arguments of this paper.

\paragraph{Organization.}
Section~\ref{sec:preliminaries} collects preliminaries, including rooted
densities.  Section~\ref{sec:row-criterion} proves the criterion for finitely
many row types.  Sections~\ref{sec:upper-tail} and~\ref{sec:fixed-densities}
apply it to prove Theorems~\ref{thm:upper-tail} and~\ref{thm:fixed-densities},
and Section~\ref{sec:conclusion} discusses open problems.
Appendix~\ref{app:arw} discusses the preprints~\cite{AlvaradoRamonWang2023,AlvaradoRamonWang2026}.

\section{Preliminaries}\label{sec:preliminaries}

This section fixes notation and collects the facts used in the proofs: basic
properties of graphons and of the two variational problems, the first
variation of homomorphism densities, the rooted densities that describe the
effect of inserting a new block of vertices, and two results from analysis in
Banach spaces.

\subsection{Graphs and graphons}\label{sec:graphons}

Throughout, $H$ is a finite simple graph with vertex set $V(H)$ and edge set
$E(H)$, and we write $N:=|V(H)|$ and $m:=|E(H)|\ge1$.  Isolated vertices do not
affect homomorphism densities, so we assume that $H$ has none.  For $v\in V(H)$,
we write $N_H(v)$ for the set of neighbors of $v$ and $d_v:=|N_H(v)|$ for its
degree, and $\Delta(H)$ for the maximum degree.  Thus $\Delta(H)=1$ exactly
when $H$ is a perfect matching.

A graphon is a symmetric measurable function $W:[0,1]^2\to[0,1]$; we identify
graphons that agree almost everywhere, and $\W_0$ denotes the space of
graphons.  Recall the edge density $e(W)$, the homomorphism density $t(H,W)$,
and the relative entropy $I_p(W)$ from~\eqref{eq:graphon-quantities}.  We use
the cut distance $\delta_\cut$ and the quotient space $\widetilde\W_0$ as
in~\cite{LovaszGraphLimits}: the space $(\widetilde\W_0,\delta_\cut)$ is
compact, $e$ and $t(H,\cdot)$ are continuous with respect to $\delta_\cut$,
and $I_p$ is lower semicontinuous~\cite{ChatterjeeVaradhanLDP}.  All objective
functions and constraints below are invariant under measure-preserving
relabelings of $[0,1]$.

Unless stated otherwise, equalities and inequalities between measurable
functions hold almost everywhere.  For functions $u_1,u_2$ on $[0,1]$, we write
$u_1\wedge u_2$ and $u_1\vee u_2$ for their pointwise minimum and maximum.  For a graphon $W$ and $x\in[0,1]$, the
function
\[
  W_x(y):=W(x,y)\qquad(y\in[0,1])
\]
is the \emph{row} of $W$ at $x$.  By Fubini's theorem, $W_x$ is measurable for
almost every $x$, and two rows are regarded as equal if they agree almost
everywhere.  A graphon is \emph{multipodal} if there are a finite measurable
partition $[0,1]=A_1\cup\dots\cup A_M$ and numbers $b_{ij}=b_{ji}\in[0,1]$ such
that $W=b_{ij}$ almost everywhere on $A_i\times A_j$ for all $i,j$; the sets
$A_i$ are its \emph{blocks}, or \emph{podes}.  Equivalently, $W$ is the
graphon of a stochastic block model with finitely many blocks.  A multipodal
graphon has finitely many distinct rows.  We say that the blocks of a
multipodal graphon are \emph{ordered} if, after relabeling the blocks,
\begin{equation}\label{eq:ordered-blocks}
  b_{ij}\le b_{i'j}\qquad\text{whenever }i\le i',\ \text{for every }j;
\end{equation}
by symmetry, the matrix $(b_{ij})$ is then nondecreasing in both indices.
Equivalently, the distinct rows $g_1,\dots,g_M$ of $W$ satisfy
$g_1\le g_2\le\dots\le g_M$ pointwise almost everywhere.

\subsection{Relative entropy and graphon entropy}

For $q\in(0,1)$ and $z\in[0,1]$, recall
$J_q(z)=z\log\frac zq+(1-z)\log\frac{1-z}{1-q}$, with $0\log0=0$.  We write
\[
  \lo(z):=\log\frac z{1-z},
  \qquad
  \sigma(\xi):=\frac1{1+e^{-\xi}}
\]
for the log-odds function and its inverse.  For $0<z<1$,
\[
  J_q'(z)=\lo(z)-\lo(q),
  \qquad
  J_q''(z)=\frac1{z(1-z)}\ge4.
\]
Thus $J_q'$ is strictly increasing from $-\infty$ to $+\infty$ on $(0,1)$ and
vanishes only at $z=q$.  Recall also the entropy density and the graphon
entropy
\[
  S_0(z)=-\frac12\bigl(z\log z+(1-z)\log(1-z)\bigr),
  \qquad
  s(W)=\int_{[0,1]^2}S_0\bigl(W(x,y)\bigr)\dd x\dd y.
\]
A direct computation gives $J_q(z)=-2S_0(z)-\log(1-q)+z\log\frac{1-q}q$, so
\begin{equation}\label{eq:entropy-identity}
  I_\eps(W)=2S_0(\eps)-2s(W)\qquad\text{whenever }e(W)=\eps.
\end{equation}
Hence, among graphons with a prescribed edge density $\eps$, minimizing
$I_\eps$ is equivalent to maximizing $s$.

\subsection{Derivatives of homomorphism densities}\label{sec:derivatives}

We regard $e$ and $t(H,\cdot)$ as polynomial functions on the Banach space
$L^\infty_{\mathrm{sym}}([0,1]^2)$ of bounded symmetric measurable functions.
For an edge $ab\in E(H)$, oriented from $a$ to $b$, and a graphon $W$, let
\[
  K_{ab,W}(x,y):=\int_{[0,1]^{V(H)\setminus\{a,b\}}}
  \prod_{ij\in E(H)\setminus\{ab\}}W(x_i,x_j)
  \prod_{i\in V(H)\setminus\{a,b\}}\dd x_i,
  \qquad x_a=x,\ x_b=y.
\]
Reversing the orientation transposes this function, so the kernel
\begin{equation}\label{eq:derivative-kernel}
  \mathcal D_W(x,y):=\frac12\sum_{ab\in E(H)}\bigl(K_{ab,W}(x,y)+K_{ab,W}(y,x)\bigr)
\end{equation}
is well defined and symmetric.  It represents the first variation of the
homomorphism density: for every bounded symmetric $U$,
\begin{equation}\label{eq:first-variation}
  Dt(H,W)[U]=\int_{[0,1]^2}\mathcal D_W(x,y)\,U(x,y)\dd x\dd y.
\end{equation}
Since each $K_{ab,W}$ takes values in $[0,1]$, we have $0\le\mathcal D_W\le m$.
If moreover $W\ge\omega>0$, then $\mathcal D_W\ge m\omega^{m-1}$.

\subsection{Rooted densities}\label{sec:rooted-densities}

The central object of this paper describes how the homomorphism density
changes when a new type of vertex is added to a graphon.  Let $W$ be a graphon
and let $u:[0,1]\to[0,1]$ be measurable.  Think of $u$ as the row of a new
vertex: $u(y)$ is the probability that the new vertex is adjacent to a vertex
with label $y$.  For $v\in V(H)$, the \emph{rooted density} of $H$ at $v$ is the
density of copies of $H$ in which the new vertex plays the role of $v$:
\begin{equation}\label{eq:rooted-density}
  \mathcal R_{v,W}(u):=\int_{[0,1]^{V(H)\setminus\{v\}}}
  \prod_{z\in N_H(v)}u(x_z)
  \prod_{\substack{ab\in E(H)\\ a,b\ne v}}W(x_a,x_b)
  \prod_{z\in V(H)\setminus\{v\}}\dd x_z,
\end{equation}
and we put
\begin{equation}\label{eq:total-rooted-density}
  \mathcal R_W(u):=\sum_{v\in V(H)}\mathcal R_{v,W}(u).
\end{equation}
If $u=W_x$ is a row of $W$, then $\mathcal R_{v,W}(W_x)$ is the homomorphism
density of $H$ in $W$ with the vertex $v$ fixed at $x$, a rooted homomorphism
density in the sense of~\cite{LovaszGraphLimits}; the definition above extends
it to rows that do not belong to $W$.  Integrating $\mathcal R_{v,W}(W_x)$ over
$x$ restores the missing variable, so
\begin{equation}\label{eq:rooted-average}
  \int_0^1\mathcal R_W(W_x)\dd x=N\,t(H,W).
\end{equation}
The map $u\mapsto\mathcal R_W(u)$ is a polynomial on $L^\infty([0,1])$ of
degree at most $\Delta(H)$, and each of its monomials has a nonnegative kernel.
Its first derivative is represented by a bounded function,
$D\mathcal R_W(u)[h]=\int_0^1\nabla\mathcal R_W(u)(y)h(y)\dd y$, where
$\nabla\mathcal R_W(u)$ is the sum over the ordered pairs $(v,z_0)$ with
$z_0\in N_H(v)$ of
\begin{equation}\label{eq:rooted-gradient}
  G_{v,z_0}(u)(y):=\int\prod_{z\in N_H(v)\setminus\{z_0\}}u(x_z)
  \prod_{\substack{ab\in E(H)\\ a,b\ne v}}W(x_a,x_b)
  \prod_{z\in V(H)\setminus\{v,z_0\}}\dd x_z,
  \qquad x_{z_0}=y.
\end{equation}
There are $2m$ such pairs and each $G_{v,z_0}(u)$ takes values in $[0,1]$, so
\begin{equation}\label{eq:rooted-lipschitz}
  0\le\nabla\mathcal R_W(u)\le2m,
  \qquad
  |\mathcal R_W(u)-\mathcal R_W(u')|\le2m\|u-u'\|_1
\end{equation}
for measurable $u,u'$ with values in $[0,1]$.  Substituting $u=W_x$ and pairing
the two orientations of each edge shows that
\begin{equation}\label{eq:rooted-gradient-derivative}
  \nabla\mathcal R_W(W_x)(y)=2\,\mathcal D_W(x,y)
  \qquad\text{for almost every }(x,y).
\end{equation}

\begin{expl}\label{expl:triangle-rooted}
For $H=K_3$, one has
$\mathcal R_W(u)=3\int_{[0,1]^2}W(y,z)u(y)u(z)\dd y\dd z$.  If $u_1$ and $u_2$ are
two rows and $h_1:=(u_1-u_2)_+$ and $h_2:=(u_2-u_1)_+$, then
\[
  \mathcal R_W(u_1\wedge u_2)+\mathcal R_W(u_1\vee u_2)
  -\mathcal R_W(u_1)-\mathcal R_W(u_2)
  =6\int_{[0,1]^2}W(y,z)\,h_1(y)h_2(z)\dd y\dd z.
\]
When $W$ is bounded below by a positive constant, the right-hand side is
positive whenever $u_1$ and $u_2$ cross, that is, whenever both
$\{u_1>u_2\}$ and $\{u_2>u_1\}$ have positive measure.  This strict
supermodularity is the mechanism behind the ordering of rows.
\end{expl}

\subsection{Two results from analysis}

We use the following two standard results.  The first is the Lagrange
multiplier rule in Banach spaces; see, for example,~\cite{Zeidler1995}.  Its
surjectivity hypothesis is the classical constraint qualification, and the
condition in Definition~\ref{def:constraint-qualification} below is a
strengthened form of it.

\begin{lemma}[Lagrange multipliers]\label{lem:lagrange}
Let $X$ be a Banach space, $\mathcal O\subseteq X$ open, and let
$f:\mathcal O\to\RR$ and $g:\mathcal O\to\RR^k$ be continuously Fr\'echet
differentiable.  If $x_0\in\mathcal O$ is a local minimizer of $f$ on
$\{x\in\mathcal O:g(x)=g(x_0)\}$ and $Dg(x_0):X\to\RR^k$ is surjective, then
there is $\lambda\in\RR^k$ with $Df(x_0)=\lambda\cdot Dg(x_0)$.
\end{lemma}

The second is the analytic implicit function theorem in Banach
spaces~\cite{Whittlesey1965}.

\begin{lemma}[Analytic implicit function theorem]\label{lem:analytic-ift}
Let $X,Y,Z$ be real Banach spaces, let $\mathcal O\subseteq X\times Y$ be
open, and let $\Xi:\mathcal O\to Z$ be real-analytic with
$\Xi(x_0,y_0)=0$.  If $D_y\Xi(x_0,y_0):Y\to Z$ is a bounded isomorphism,
then there are neighborhoods $\mathcal U\ni x_0$ and $\mathcal V\ni y_0$ and a
real-analytic map $\varphi:\mathcal U\to\mathcal V$ such that, for
$(x,y)\in\mathcal U\times\mathcal V$, $\Xi(x,y)=0$ if and only if
$y=\varphi(x)$.
\end{lemma}

\section{A criterion for finitely many row types}\label{sec:row-criterion}

This section contains the common core of the proofs of
Theorems~\ref{thm:upper-tail} and~\ref{thm:fixed-densities}.  It is stated
without reference to either variational problem.  The input is a functional on
rows whose global minimum is attained by the rows of a graphon $W$; the output
is that $W$ is multipodal with ordered blocks.

Throughout this section, we fix a graph $H$ with $\Delta(H)\ge2$, a graphon
$W$, and numbers $q\in(0,1)$, $\alpha\in\RR$, $\mu>0$, and $\omega>0$ such that
\begin{equation}\label{eq:positive-base}
  W\ge\omega\qquad\text{almost everywhere}.
\end{equation}
For measurable $u:[0,1]\to[0,1]$, define the \emph{row functional}
\begin{equation}\label{eq:row-functional}
  \mathcal F_W(u):=2\int_0^1J_q\bigl(u(y)\bigr)\dd y-2\alpha\int_0^1u(y)\dd y
  -\mu\,\mathcal R_W(u).
\end{equation}
Assume that there is $c_W\in\RR$ such that
\begin{equation}\label{eq:row-lower-bound}
  \mathcal F_W(u)\ge c_W\qquad\text{for every measurable }u:[0,1]\to[0,1],
\end{equation}
and let
\[
  \mathcal M_W:=\bigl\{u\in L^\infty([0,1]):0\le u\le1,\ \mathcal F_W(u)=c_W\bigr\}
\]
be the set of global minimizers of $\mathcal F_W$ attaining this bound.

\begin{theorem}[Finitely many row types]\label{thm:row-criterion}
Under the assumptions above, the set $\mathcal M_W$ is finite, and any two of
its elements $u_1,u_2$ satisfy $u_1\le u_2$ or $u_2\le u_1$ almost everywhere.
If, in addition,
\begin{equation}\label{eq:rows-attain}
  W_x\in\mathcal M_W\qquad\text{for almost every }x\in[0,1],
\end{equation}
then $W$ is multipodal and its blocks are ordered.
\end{theorem}

The proof has three steps.  First, the singularities of $J_q'$ at $0$ and $1$
confine all minimizers to a common interval $[\delta,1-\delta]$, where they
satisfy an Euler--Lagrange equation.  Second, strict supermodularity of
$\mathcal R_W$ shows that two minimizers cannot cross, so $\mathcal M_W$ is
totally ordered.  Third, compactness and an analysis of the second variation
of $\mathcal F_W$ show that a totally ordered set of minimizers cannot contain
a continuum, and is therefore finite.  The assumption~\eqref{eq:positive-base}
is used only in the strict inequalities of the second and third steps.

\subsection{Interiority and the Euler--Lagrange equation}

\begin{lemma}[Interiority]\label{lem:row-interior}
There is $\delta\in(0,1/2)$ such that every $u\in\mathcal M_W$ satisfies
$\delta\le u\le1-\delta$.  Moreover, every $u\in\mathcal M_W$ satisfies
\begin{equation}\label{eq:row-euler}
  2J_q'(u)-2\alpha=\mu\nabla\mathcal R_W(u),
  \qquad\text{equivalently}\qquad
  u=\sigma\Bigl(\lo(q)+\alpha+\frac\mu2\nabla\mathcal R_W(u)\Bigr).
\end{equation}
\end{lemma}

\begin{proof}
Choose $\delta\in(0,1/2)$ so small that $2J_q'(z)-2\alpha+2m\mu<0$ for
$0<z<\delta$ and $2J_q'(z)-2\alpha-2m\mu>0$ for $1-\delta<z<1$.  Let
$u\in\mathcal M_W$ and suppose that $A:=\{u<\delta\}$ has positive measure.
Put $\tilde u:=\max(u,\delta)$.  By~\eqref{eq:rooted-lipschitz},
\[
  \mathcal F_W(\tilde u)-\mathcal F_W(u)
  \le\int_A\Bigl(2\bigl(J_q(\delta)-J_q(u)\bigr)-2\alpha(\delta-u)
  +2m\mu(\delta-u)\Bigr)
  =\int_A\int_u^\delta\bigl(2J_q'(z)-2\alpha+2m\mu\bigr)\dd z<0,
\]
contradicting~\eqref{eq:row-lower-bound}.  The case $\{u>1-\delta\}$ is
analogous, with $\tilde u:=\min(u,1-\delta)$.  Since every $u\in\mathcal M_W$
takes values in $[\delta,1-\delta]$, the functions $u+th$ with bounded $h$ and
small $|t|$ are admissible, so the first variation of $\mathcal F_W$ at $u$
vanishes in every bounded direction.  This gives~\eqref{eq:row-euler}.
\end{proof}

\subsection{Minimizers cannot cross}

\begin{lemma}[Strict supermodularity]\label{lem:strict-supermodular}
Let $u_1,u_2\in\mathcal M_W$.  Then
\[
  \mathcal R_W(u_1\wedge u_2)+\mathcal R_W(u_1\vee u_2)
  \ge\mathcal R_W(u_1)+\mathcal R_W(u_2),
\]
and the inequality is strict if $u_1$ and $u_2$ cross, that is, if both
$\{u_1>u_2\}$ and $\{u_2>u_1\}$ have positive measure.
\end{lemma}

\begin{proof}
Put $g:=u_1\wedge u_2$, $h_1:=(u_1-u_2)_+$, and $h_2:=(u_2-u_1)_+$, so that
$u_1=g+h_1$, $u_2=g+h_2$, and $u_1\vee u_2=g+h_1+h_2$.  Expanding the products
in~\eqref{eq:rooted-density}, the difference of the two sides is a sum of
integrals of nonnegative kernels against products containing at least one
factor $h_1$ and at least one factor $h_2$.  This proves the weak inequality.
For strictness, choose a vertex $v_0$ with degree $d_0\ge2$.  Among the terms
of $\mathcal R_{v_0,W}$, keep only those with exactly one factor $h_1$, exactly
one factor $h_2$, and $d_0-2$ factors $g$.  Each edge factor not incident to
$v_0$ is at least $\omega$ by~\eqref{eq:positive-base}, and $g\ge\delta$ by
Lemma~\ref{lem:row-interior}, so
\[
  \mathcal R_W(u_1\wedge u_2)+\mathcal R_W(u_1\vee u_2)
  -\mathcal R_W(u_1)-\mathcal R_W(u_2)
  \ge d_0(d_0-1)\,\omega^{m-d_0}\delta^{d_0-2}
  \int_0^1h_1\int_0^1h_2,
\]
which is positive when $u_1$ and $u_2$ cross.
\end{proof}

\begin{prop}[Total order]\label{prop:row-chain}
For all $u_1,u_2\in\mathcal M_W$, either $u_1\le u_2$ or $u_2\le u_1$ almost
everywhere.
\end{prop}

\begin{proof}
Pointwise, the two numbers $(u_1\wedge u_2)(y)$ and $(u_1\vee u_2)(y)$ are
$u_1(y)$ and $u_2(y)$ in some order.  Hence the sum of the first two terms
of~\eqref{eq:row-functional} over the pair $u_1\wedge u_2$, $u_1\vee u_2$ equals
the corresponding sum over the pair $u_1$, $u_2$.  If $u_1$ and $u_2$ crossed,
Lemma~\ref{lem:strict-supermodular} and $\mu>0$ would give
\[
  \mathcal F_W(u_1\wedge u_2)+\mathcal F_W(u_1\vee u_2)
  <\mathcal F_W(u_1)+\mathcal F_W(u_2)=2c_W,
\]
contradicting~\eqref{eq:row-lower-bound} applied to $u_1\wedge u_2$ and
$u_1\vee u_2$.
\end{proof}

\subsection{Compactness and the second variation}

\begin{lemma}[Compactness]\label{lem:row-compact}
The set $\mathcal M_W$ is compact in $L^\infty([0,1])$.
\end{lemma}

\begin{proof}
Each term~\eqref{eq:rooted-gradient} of $\nabla\mathcal R_W(u)$ has the form
$y\mapsto\int K(y,z_1,\dots,z_k)\prod_{i=1}^ku(z_i)\dd z$ with a measurable
kernel $0\le K\le1$.  Such a kernel defines a Hilbert--Schmidt, hence compact,
operator from $L^2([0,1]^k)$ to $L^2([0,1])$, and the functions
$u^{\otimes k}$ with $0\le u\le1$ form a bounded subset of $L^2([0,1]^k)$.
Hence $\{\nabla\mathcal R_W(u):0\le u\le1\}$ is relatively compact in $L^2$.
Moreover, expanding the difference of products telescopically gives
\begin{equation}\label{eq:gradient-lipschitz}
  \|\nabla\mathcal R_W(u)-\nabla\mathcal R_W(u')\|_\infty\le C_H\|u-u'\|_1.
\end{equation}
Now let $(u_j)\subseteq\mathcal M_W$.  After passing to a subsequence,
$\nabla\mathcal R_W(u_j)\to g$ in $L^2$.  Since $\sigma$ is Lipschitz, the
fixed-point form of~\eqref{eq:row-euler} gives
$u_j\to u:=\sigma(\lo(q)+\alpha+\frac\mu2g)$ in $L^2$.
By~\eqref{eq:gradient-lipschitz}, $\nabla\mathcal R_W(u_j)\to\nabla\mathcal R_W(u)$
uniformly, so $g=\nabla\mathcal R_W(u)$, and~\eqref{eq:row-euler} now gives
$u_j\to u$ uniformly.  Since $\mathcal F_W$ is continuous under uniform
convergence, $u\in\mathcal M_W$.
\end{proof}

For $u\in\mathcal M_W$, consider the second variation
\[
  Q_u(h):=2\int_0^1J_q''(u)\,h^2-\mu\,D^2\mathcal R_W(u)[h,h],
  \qquad h\in L^2([0,1]).
\]
The bilinear form $D^2\mathcal R_W(u)$ is given by a bounded symmetric kernel,
so it is represented by a compact self-adjoint operator $B_u$ on
$L^2([0,1])$.  By Lemma~\ref{lem:row-interior}, multiplication by
$2J_q''(u)=2/(u(1-u))$ is a bounded operator $A_u$ with $A_u\ge8$ and bounded
inverse.  Thus $Q_u(h)=\langle h,\mathcal L_uh\rangle$ with
$\mathcal L_u:=A_u-\mu B_u$.

\begin{lemma}[Coercivity on mean-zero functions]\label{lem:coercive}
For every $u\in\mathcal M_W$, there is $c_u>0$ such that
$Q_u(h)\ge c_u\|h\|_2^2$ for every $h\in L^2([0,1])$ with $\int_0^1h=0$.
\end{lemma}

\begin{proof}
Since $u$ is a global minimizer and takes values in $[\delta,1-\delta]$, the
second variation is nonnegative: $Q_u(h)\ge0$ for bounded $h$, and by density
for all $h\in L^2$.  Let $h\ne0$ have mean zero, and write $h=h_+-h_-$ with
$h_\pm\ge0$ of disjoint supports; both are nonzero.  Choosing a vertex $v_0$ of
degree $d_0\ge2$ as in the proof of Lemma~\ref{lem:strict-supermodular}, we get
$D^2\mathcal R_W(u)[h_+,h_-]\ge d_0(d_0-1)\omega^{m-d_0}\delta^{d_0-2}
\|h_+\|_1\|h_-\|_1>0$.  Since $h_+$ and $h_-$ have disjoint supports,
\[
  Q_u(h)=Q_u(h_+)+Q_u(h_-)+2\mu\,D^2\mathcal R_W(u)[h_+,h_-]>0.
\]
Suppose that the quantitative bound fails.  Then there are $h_j$ with mean
zero, $\|h_j\|_2=1$, and $Q_u(h_j)\to0$.  Passing to a subsequence,
$h_j\rightharpoonup h$ weakly in $L^2$, and $h$ has mean zero.  Compactness of
$B_u$ gives $\langle h_j,B_uh_j\rangle\to\langle h,B_uh\rangle$.  If $h=0$, then
$Q_u(h_j)\ge8-\mu\langle h_j,B_uh_j\rangle\to8$, a contradiction.  If $h\ne0$,
then weak lower semicontinuity of $h\mapsto\langle h,A_uh\rangle$ gives
$Q_u(h)\le\liminf Q_u(h_j)=0$, contradicting strict positivity.
\end{proof}

\begin{lemma}[Bordered second variation]\label{lem:bordered}
For every $u\in\mathcal M_W$, the operator
$(h,\lambda)\mapsto\bigl(\mathcal L_uh-\lambda\mathbf 1,\int_0^1h\bigr)$ is a
bounded isomorphism of $L^2([0,1])\times\RR$ onto itself and of
$L^\infty([0,1])\times\RR$ onto itself.
\end{lemma}

\begin{proof}
Let $L^2_0$ be the space of mean-zero functions and $P_0$ the orthogonal
projection onto it.  By Lemma~\ref{lem:coercive} and the Lax--Milgram theorem,
$P_0\mathcal L_u$ restricted to $L^2_0$ is invertible.  Given
$(g,\beta)\in L^2\times\RR$, write $h=\beta\mathbf1+k$ with $k\in L^2_0$, solve
$P_0\mathcal L_uk=P_0(g-\mathcal L_u(\beta\mathbf 1))$, and put
$\lambda:=\int_0^1(\mathcal L_uh-g)$; then $\mathcal L_uh-\lambda\mathbf1=g$.  If
$(h,\lambda)$ is in the kernel, then $h\in L^2_0$ and
$Q_u(h)=\lambda\int h=0$, so $h=0$ and $\lambda=0$.  This proves the claim on
$L^2\times\RR$.  If $g\in L^\infty$, then
$h=A_u^{-1}(g+\lambda\mathbf1+\mu B_uh)$, and $B_u$ maps $L^2$ boundedly into
$L^\infty$ because its kernel is bounded; hence $h\in L^\infty$ with a
corresponding norm bound.
\end{proof}

\subsection{Finiteness}

\begin{lemma}[Finiteness]\label{lem:row-finite}
The set $\mathcal M_W$ is finite.
\end{lemma}

\begin{proof}
We may assume $\mathcal M_W\ne\emptyset$.  By Proposition~\ref{prop:row-chain},
distinct elements of $\mathcal M_W$ differ on a set of positive measure with a
fixed sign, so they have distinct means.  Hence the map
$\mathfrak m(u):=\int_0^1u$ is injective on $\mathcal M_W$, and by
Lemma~\ref{lem:row-compact} it is a homeomorphism from $\mathcal M_W$ onto the
compact set $\mathcal A:=\mathfrak m(\mathcal M_W)\subseteq\RR$.

Fix $u_*\in\mathcal M_W$ and put $a_*:=\mathfrak m(u_*)$.  Consider the map
\[
  \Xi(a;u,\lambda):=\Bigl(2J_q'(u)-2\alpha-\mu\nabla\mathcal R_W(u)-\lambda\mathbf1,\
  \int_0^1u-a\Bigr),
\]
defined for $a\in\RR$, $\lambda\in\RR$, and $u$ in the open $L^\infty$-ball of
radius $\delta/2$ around $u_*$.  It is real-analytic with values in
$L^\infty\times\RR$: the map $u\mapsto J_q'(u)$ is analytic on this ball, and
$u\mapsto\nabla\mathcal R_W(u)$ is polynomial.  We have $\Xi(a_*;u_*,0)=0$ by
\eqref{eq:row-euler}, and the derivative of $\Xi$ in $(u,\lambda)$ is the
isomorphism of Lemma~\ref{lem:bordered}.  By Lemma~\ref{lem:analytic-ift},
there are an open interval $I_*\ni a_*$, a neighborhood of $(u_*,0)$, and
analytic maps $a\mapsto u_*(a)$ and $a\mapsto\lambda_*(a)$ on $I_*$ such that,
in this neighborhood, $\Xi(a;u,\lambda)=0$ if and only if
$(u,\lambda)=(u_*(a),\lambda_*(a))$.

Suppose that $a_*$ is an accumulation point of $\mathcal A$, and choose distinct
$a_j\in\mathcal A$ with $a_j\to a_*$ and $u_j\in\mathcal M_W$ with
$\mathfrak m(u_j)=a_j$.  Then $u_j\to u_*$ in $L^\infty$, and each $u_j$ satisfies
$\Xi(a_j;u_j,0)=0$ by~\eqref{eq:row-euler}.  By local uniqueness,
$\lambda_*(a_j)=0$ for all large $j$.  The analytic function $\lambda_*$ has
zeros accumulating at $a_*$, so $\lambda_*\equiv0$ near $a_*$.  Along the
branch, differentiating $\int u_*(a)=a$ gives
\[
  \frac{\dd}{\dd a}\mathcal F_W\bigl(u_*(a)\bigr)
  =\int_0^1\bigl(2J_q'(u_*(a))-2\alpha-\mu\nabla\mathcal R_W(u_*(a))\bigr)u_*'(a)
  =\lambda_*(a)=0.
\]
Hence $\mathcal F_W(u_*(a))=c_W$ near $a_*$, so $u_*(a)\in\mathcal M_W$ and a
whole open interval around $a_*$ lies in $\mathcal A$.

We have shown that every accumulation point of $\mathcal A$ is an interior
point of $\mathcal A$.  If $\mathcal A$ were infinite, its set $\mathcal A'$ of
accumulation points would be nonempty and compact.  By what we have shown,
every point of $\mathcal A'$ has a neighborhood contained in $\mathcal A$, all
of whose points are accumulation points; hence $\mathcal A'$ is open.  A
nonempty subset of $\RR$ cannot be both open and compact, so $\mathcal A$, and
therefore $\mathcal M_W$, is finite.
\end{proof}

\begin{proof}[Proof of Theorem~\ref{thm:row-criterion}]
Proposition~\ref{prop:row-chain} and Lemma~\ref{lem:row-finite} show that
$\mathcal M_W=\{g_1,\dots,g_M\}$ is a finite totally ordered set; label its
elements so that $g_1\le\dots\le g_M$.  Assume~\eqref{eq:rows-attain}, and let
$A_i:=\{x\in[0,1]:W_x=g_i\text{ almost everywhere}\}$.  These sets are
measurable, pairwise disjoint, and cover $[0,1]$ up to a null set; we discard
the null sets among them and add the remaining null set to one of the others.
For almost every $(x,y)\in A_i\times A_j$, the row identities and the symmetry
of $W$ give $g_i(y)=W(x,y)=W(y,x)=g_j(x)$.  A function of $y$ that agrees almost
everywhere on $A_i\times A_j$ with a function of $x$ is almost everywhere
constant, so $W=b_{ij}$ almost everywhere on $A_i\times A_j$ for some constant
$b_{ij}=b_{ji}$.  Thus $W$ is multipodal.  Finally, $g_i=b_{ij}$ on $A_j$, so
$g_i\le g_{i'}$ for $i\le i'$ is exactly~\eqref{eq:ordered-blocks}.
\end{proof}

\begin{rmk}\label{rmk:criterion-input}
Theorem~\ref{thm:row-criterion} does not use that $W$ solves any graphon
variational problem, nor an Euler--Lagrange equation for $W$ itself.  All
information about the ambient problem enters through the global lower
bound~\eqref{eq:row-lower-bound}, the sign $\mu>0$, the lower
bound~\eqref{eq:positive-base}, and the property~\eqref{eq:rows-attain}.  In
Sections~\ref{sec:upper-tail} and~\ref{sec:fixed-densities}, these hypotheses
are supplied by inserting a small block of vertices into an optimizer.
\end{rmk}

\section{The upper-tail problem}\label{sec:upper-tail}

In this section we prove Theorem~\ref{thm:upper-tail}.  Fix $0<p<r<1$.  The
infimum in~\eqref{eq:graphon-variational-problem} is attained, because the
feasible set is closed in the compact space $(\widetilde\W_0,\delta_\cut)$ and
$I_p$ is lower semicontinuous.  If $H$ is a perfect matching, then
$t(H,W)=e(W)^m$, and feasibility means $e(W)\ge r$; Jensen's inequality and the
monotonicity of $J_p$ on $[r,1]$ give $I_p(W)\ge J_p(e(W))\ge J_p(r)$, with
equality only for $W\equiv r$.  Thus $W\equiv r$ is the unique minimizer, and
it is trivially multipodal.  For the rest of this section we assume
$\Delta(H)\ge2$, and we verify the hypotheses of
Theorem~\ref{thm:row-criterion} for an arbitrary minimizer $W$.

\subsection{Basic properties of minimizers}

\begin{lemma}\label{lem:upper-tail-basic}
Let $W$ be a minimizer of~\eqref{eq:graphon-variational-problem}.  Then
$W\ge p$ almost everywhere, the constraint is active, $t(H,W)=r^m$, and there is
$\gamma>0$ such that $W\le1-\gamma$ almost everywhere.
\end{lemma}

\begin{proof}
\emph{Lower bound.}  Put $W^+:=\max(W,p)$.  Since the factors in $t(H,\cdot)$ are
nonnegative, $t(H,W^+)\ge t(H,W)$, so $W^+$ is feasible.  Since $J_p$ is strictly
decreasing on $[0,p]$, we have $I_p(W^+)\le I_p(W)$, with strict inequality
unless $W\ge p$ almost everywhere.

\emph{Active constraint.}  Suppose that $t(H,W)>r^m$, and put
$W_\lambda:=p+\lambda(W-p)$ for $\lambda\in[0,1]$.  The function
$\lambda\mapsto t(H,W_\lambda)$ is continuous, and it equals $p^m<r^m$ at
$\lambda=0$ and exceeds $r^m$ at $\lambda=1$.  Hence $t(H,W_{\lambda_0})=r^m$
for some $\lambda_0\in(0,1)$.  Since $J_p$ is strictly convex with $J_p(p)=0$
and $W\not\equiv p$, we get $I_p(W_{\lambda_0})<\lambda_0I_p(W)\le I_p(W)$, a
contradiction.

\emph{Upper bound.}  Since $W\ge p$, the derivative kernel satisfies
$mp^{m-1}\le\mathcal D_W\le m$.  Because $t(H,W)=r^m<1$, there are $\gamma_0>0$
and a symmetric set $A$ of positive measure with $W\le1-\gamma_0$ on $A$.
Suppose that $W$ is not essentially bounded away from $1$.  For large
$\Lambda>0$, the symmetric set $E_\Lambda:=\{W>1-e^{-\Lambda}\}$ has positive
measure and is disjoint from $A$.  Put
\[
  c_\Lambda:=\frac{\int_{E_\Lambda}\mathcal D_W}{\int_A\mathcal D_W},
  \qquad
  U_\Lambda:=-\mathbf 1_{E_\Lambda}+c_\Lambda\mathbf 1_A,
\]
so that $Dt(H,W)[U_\Lambda]=0$ by~\eqref{eq:first-variation} and
$c_\Lambda|A|\le p^{1-m}|E_\Lambda|$.  The function
$(\theta,\xi)\mapsto t(H,W+\theta U_\Lambda+\xi\mathbf 1_A)$ is a polynomial, and
its derivative in $\xi$ at the origin is $\int_A\mathcal D_W>0$.  The implicit
function theorem gives $\xi_\Lambda(\theta)=O_\Lambda(\theta^2)$ such that
$V_\theta:=W+\theta U_\Lambda+\xi_\Lambda(\theta)\mathbf 1_A$ satisfies
$t(H,V_\theta)=t(H,W)$ for small $\theta$.  For $0<\theta<e^{-\Lambda}/2$ small
enough, $V_\theta$ is a graphon: it equals $W-\theta\ge p$ on $E_\Lambda$, and
on $A$ it equals $W+\theta c_\Lambda+\xi_\Lambda(\theta)$, where
$0\le\theta c_\Lambda+\xi_\Lambda(\theta)\le\gamma_0/2$ because $c_\Lambda>0$
and $\xi_\Lambda(\theta)=O_\Lambda(\theta^2)$.  In particular, the values of
$V_\theta$ on $A$ stay in $[p,1-\gamma_0/2]$, where $|J_p'|$ is bounded by a
constant $C$.  If $z\in[1-e^{-\Lambda},1]$, every point between $z-\theta$ and
$z$ is at least $1-\frac32e^{-\Lambda}$, where $J_p'\ge\Lambda-C_p$ for a
constant $C_p$ independent of $\Lambda$; hence
$J_p(z-\theta)-J_p(z)\le-\theta(\Lambda-C_p)$.  Therefore
\[
  I_p(V_\theta)-I_p(W)
  \le-\theta(\Lambda-C_p)|E_\Lambda|+C\theta c_\Lambda|A|+O_\Lambda(\theta^2)
  \le-\theta\bigl(\Lambda-C_p-Cp^{1-m}\bigr)|E_\Lambda|+O_\Lambda(\theta^2).
\]
For $\Lambda$ large and then $\theta>0$ small, this is negative, contradicting
the minimality of $W$.
\end{proof}

\begin{lemma}[Positive multiplier]\label{lem:upper-tail-multiplier}
Let $W$ be a minimizer of~\eqref{eq:graphon-variational-problem}.  There is
$\mu>0$ such that
\begin{equation}\label{eq:upper-tail-el}
  J_p'\bigl(W(x,y)\bigr)=\mu\,\mathcal D_W(x,y)\qquad\text{for almost every }(x,y).
\end{equation}
In particular, $W$ takes values in a compact subinterval of $(p,1)$.
\end{lemma}

\begin{proof}
By Lemma~\ref{lem:upper-tail-basic}, there is $\gamma\in(0,p]$ with
$\gamma\le W\le1-\gamma$.  On the open ball of radius $\gamma/2$ around $W$ in
$L^\infty_{\mathrm{sym}}([0,1]^2)$, every element is a graphon, $I_p$ is
continuously differentiable, and $t(H,\cdot)$ is polynomial.  Since the
constraint is active, $W$ minimizes $I_p$ on the intersection of this ball with
$\{t(H,\cdot)=r^m\}$, and $Dt(H,W)[\mathbf1]=\int\mathcal D_W>0$.  By
Lemma~\ref{lem:lagrange}, $DI_p(W)=\mu\,Dt(H,W)$ for some $\mu\in\RR$.  The
bounded symmetric function $g:=J_p'(W)-\mu\mathcal D_W$ then satisfies
$\int gU=0$ for every bounded symmetric $U$, and taking $U=\operatorname{sign}g$
shows $g=0$, which is~\eqref{eq:upper-tail-el}.  Since $W+\theta\mathbf 1$ is
feasible for small $\theta>0$, we have $DI_p(W)[\mathbf 1]\ge0$, and hence
$\mu\ge0$.  If $\mu=0$, then $J_p'(W)=0$ and $W\equiv p$, which is not feasible.
So $\mu>0$, and~\eqref{eq:upper-tail-el} with $mp^{m-1}\le\mathcal D_W\le m$
gives $\sigma(\lo(p)+\mu mp^{m-1})\le W\le\sigma(\lo(p)+\mu m)$.
\end{proof}

\subsection{Inserting a block of vertices}

The next lemma is the key step.  It upgrades the Euler--Lagrange
equation~\eqref{eq:upper-tail-el}, which only says that the rows of $W$ are
critical points of a row functional, to the statement that they are global
minimizers of it.

\begin{lemma}[Insertion inequality]\label{lem:upper-tail-insertion}
Let $W$ be a minimizer of~\eqref{eq:graphon-variational-problem} and let
$\mu>0$ be as in Lemma~\ref{lem:upper-tail-multiplier}.  Let $\mathcal F_W$ be
the row functional~\eqref{eq:row-functional} with $q=p$ and $\alpha=0$.  Then
\[
  \mathcal F_W(u)\ge c_W:=2I_p(W)-\mu N\,t(H,W)
  \qquad\text{for every measurable }u:[0,1]\to[0,1],
\]
and $\mathcal F_W(W_x)=c_W$ for almost every $x\in[0,1]$.
\end{lemma}

\begin{proof}
Fix a measurable $u:[0,1]\to[0,1]$ and $\theta\in(0,1)$.  Let
$\phi_\theta:[\theta,1]\to[0,1]$ be the increasing affine bijection.  For
$\xi\in\RR$ with $|\xi|$ small, define the symmetric function
$W_{\theta,\xi}$ on $[0,1]^2$ by
\[
  W_{\theta,\xi}(x,y):=
  \begin{cases}
    p, & x,y<\theta,\\
    u(\phi_\theta(y)), & x<\theta\le y,\\
    u(\phi_\theta(x)), & y<\theta\le x,\\
    W(\phi_\theta(x),\phi_\theta(y))+\xi, & x,y\ge\theta.
  \end{cases}
\]
In words, we insert a new block of vertices of measure $\theta$ whose row is $u$,
keep $W$ on the old vertices, and shift the old values by $\xi$.  By
Lemma~\ref{lem:upper-tail-basic}, $W_{\theta,\xi}$ is a graphon for $|\xi|$ small.
Splitting the integral defining the homomorphism density according to the set
$S\subseteq V(H)$ of vertices whose variables lie in the new block gives
\[
  P(\theta,\xi):=t(H,W_{\theta,\xi})
  =\sum_{S\subseteq V(H)}\theta^{|S|}(1-\theta)^{N-|S|}T_S(\xi),
\]
where each $T_S$ is a polynomial in $\xi$ that does not depend on $\theta$.
Here $T_\emptyset(\xi)=t(H,W+\xi)$ and $T_{\{v\}}(0)=\mathcal R_{v,W}(u)$.
The right-hand side is a polynomial in $(\theta,\xi)$.  We have
$P(0,0)=t(H,W)$ and $\partial_\xi P(0,0)=Dt(H,W)[\mathbf1]>0$, so the implicit
function theorem gives a smooth function $\xi(\theta)$ with $\xi(0)=0$
and $P(\theta,\xi(\theta))=t(H,W)$.  Differentiating at $\theta=0$ gives
\[
  \xi'(0)\,Dt(H,W)[\mathbf1]=-\partial_\theta P(0,0)
  =-\bigl(\mathcal R_W(u)-N\,t(H,W)\bigr).
\]
For small $\theta>0$, the graphon $W_\theta:=W_{\theta,\xi(\theta)}$ is feasible, so
$I_p(W_\theta)\ge I_p(W)$.  The relative entropy has the exact expression
\[
  I_p(W_\theta)=(1-\theta)^2\int_{[0,1]^2}J_p\bigl(W+\xi(\theta)\bigr)
  +2\theta(1-\theta)\int_0^1J_p(u)+\theta^2J_p(p),
\]
which is differentiable at $\theta=0$.  Using
$DI_p(W)[\mathbf1]=\mu\,Dt(H,W)[\mathbf 1]$ from
Lemma~\ref{lem:upper-tail-multiplier}, we obtain
\[
  0\le\frac{\dd}{\dd\theta}\Big|_{\theta=0^+}I_p(W_\theta)
  =-2I_p(W)+2\int_0^1J_p(u)-\mu\bigl(\mathcal R_W(u)-N\,t(H,W)\bigr)
  =\mathcal F_W(u)-c_W.
\]
This proves the inequality.  For the rows, Fubini's theorem gives
$\int_0^1\int_0^1J_p(W_x(y))\dd y\dd x=I_p(W)$, and~\eqref{eq:rooted-average}
gives $\int_0^1\mathcal R_W(W_x)\dd x=N\,t(H,W)$.  Hence
$\int_0^1\mathcal F_W(W_x)\dd x=c_W$, and since $\mathcal F_W(W_x)\ge c_W$ for
almost every $x$, equality holds for almost every $x$.
\end{proof}

\begin{proof}[Proof of Theorem~\ref{thm:upper-tail}]
The case of a perfect matching was treated at the beginning of this section.
If $\Delta(H)\ge2$ and $W$ is a minimizer, apply Theorem~\ref{thm:row-criterion}
with $q=p$, $\alpha=0$, $\omega=p$, and $\mu>0$ from
Lemma~\ref{lem:upper-tail-multiplier}.  The lower bound~\eqref{eq:positive-base}
holds by Lemma~\ref{lem:upper-tail-basic}, and Lemma~\ref{lem:upper-tail-insertion}
provides~\eqref{eq:row-lower-bound} and~\eqref{eq:rows-attain}.
\end{proof}

\begin{rmk}[Why global optimality is needed]\label{rmk:why-global}
The Euler--Lagrange equation~\eqref{eq:upper-tail-el} alone implies only that
almost every row of $W$ is a critical point of $\mathcal F_W$.  For stars, the
equation has a special structure: if $H=K_{1,k}$, then
$\mathcal D_W(x,y)=\frac k2\bigl(d_W(x)^{k-1}+d_W(y)^{k-1}\bigr)$ with
$d_W(x)=\int_0^1W(x,y)\dd y$, so~\eqref{eq:upper-tail-el} expresses $W(x,y)$
through a sum of a function of $x$ and a function of $y$, and much of the
analysis reduces to a one-variable equation for $d_W$.  A separation of this
kind underlies the results of Kenyon, Radin, Ren, and
Sadun~\cite{KenyonRadinRenSadun2017} for edge and star constraints.  For general $H$ no such separation is available;
already for the $4$-cycle, $\mathcal D_W$ is a multiple of the threefold
operator product $\int W(x,z_1)W(z_1,z_2)W(z_2,y)\dd z_1\dd z_2$.  The insertion
inequality of Lemma~\ref{lem:upper-tail-insertion} is what makes the argument
work for all graphs.  It cannot be replaced by second-order conditions at $W$:
inserting a small block is a small perturbation in $L^1$ but not in
$L^\infty$, and it changes the value of $I_p$ at first order in the measure of
the block; see Appendix~\ref{app:arw-second-order}.
\end{rmk}

\section{Prescribed edge and subgraph densities}\label{sec:fixed-densities}

In this section we prove Theorem~\ref{thm:fixed-densities} and
Corollary~\ref{cor:almost-every-tau}.  Fix a graph $H$ with $\Delta(H)\ge2$
and $\eps\in(0,1)$, and recall
$\tau_{\max}(\eps)=\max\{t(H,W):e(W)=\eps\}$.  By~\eqref{eq:entropy-identity},
the maximizers of~\eqref{eq:entropy-maximization} are exactly the minimizers of
\begin{equation}\label{eq:fixed-problem}
  \Psi_\eps(\tau):=\min\bigl\{I_\eps(W):e(W)=\eps,\ t(H,W)=\tau\bigr\},
  \qquad
  \eps^m<\tau<\tau_{\max}(\eps),
\end{equation}
and we work with this formulation.  The interval of values of $\tau$ is
nonempty: the graphon $\mathbf 1_{[0,\sqrt\eps]^2}$ has edge density $\eps$ and
$H$-density $\eps^{N/2}$, and $\eps^{N/2}>\eps^m$ because $2m\ge N+1$ when
$\Delta(H)\ge2$.  Every value in the interval is attained along the segment
$\eps+\theta(W_{\max}-\eps)$, $\theta\in[0,1]$, where $W_{\max}$ attains
$\tau_{\max}(\eps)$, and minimizers exist by compactness.

\subsection{The constraint qualification}

Inserting a block of vertices changes both the edge density and the
$H$-density, and both must be restored.  This requires two independent
directions of perturbation.

\begin{defn}[Constraint qualification]\label{def:constraint-qualification}
Let $W$ be a graphon with $e(W)=\eps$ and $t(H,W)=\tau$.  For $\eta>0$, put
$\Omega_\eta(W):=\{(x,y):\eta<W(x,y)<1-\eta\}$.  We say that $W$
\emph{satisfies the constraint qualification} if there are $\eta>0$ and bounded
symmetric functions $U_1,U_2$ supported in $\Omega_\eta(W)$ such that the matrix
\begin{equation}\label{eq:cq-matrix}
  \Gamma_W(U_1,U_2):=
  \begin{pmatrix}
    \int U_1 & \int U_2\\[1mm]
    \int\mathcal D_WU_1 & \int\mathcal D_WU_2
  \end{pmatrix}
\end{equation}
is invertible.  Otherwise we call $W$ \emph{degenerate}.
\end{defn}

The rows of $\Gamma_W(U_1,U_2)$ are the first variations of $e$ and
$t(H,\cdot)$ at $W$ in the directions $U_1$ and $U_2$.  Thus the constraint
qualification says that the derivative of $V\mapsto(e(V),t(H,V))$ at $W$ is
surjective onto $\RR^2$, as required in Lemma~\ref{lem:lagrange}, even when
restricted to bounded perturbations supported where $W$ is bounded away from
$0$ and $1$.  The support condition ensures that small corrections along
$U_1,U_2$ keep the graphon inside $[0,1]$.

\subsection{Monotonicity and positive multipliers}

\begin{lemma}[Monotonicity]\label{lem:fixed-monotone}
The function $\tau\mapsto\Psi_\eps(\tau)$ is strictly increasing on
$(\eps^m,\tau_{\max}(\eps))$.  Moreover, a minimizer of~\eqref{eq:fixed-problem}
at $\tau$ also minimizes $I_\eps$ over all graphons $U$ with $e(U)=\eps$ and
$t(H,U)\ge\tau$.
\end{lemma}

\begin{proof}
Let $\eps^m<\tau_1<\tau_2<\tau_{\max}(\eps)$, and let $W_2$ be a minimizer at
$\tau_2$.  The graphons $W_\theta:=\eps+\theta(W_2-\eps)$ have edge density
$\eps$, and $\theta\mapsto t(H,W_\theta)$ is continuous with values $\eps^m$ and
$\tau_2$ at $\theta=0,1$; hence $t(H,W_{\theta_0})=\tau_1$ for some
$\theta_0\in(0,1)$.  By strict convexity of $J_\eps$ and $J_\eps(\eps)=0$,
$\Psi_\eps(\tau_1)\le I_\eps(W_{\theta_0})<\theta_0I_\eps(W_2)\le\Psi_\eps(\tau_2)$.
For the second claim, if $e(U)=\eps$ and $t(H,U)>\tau$, the same interpolation
between $\eps$ and $U$ produces a graphon with $H$-density $\tau$ and cost
strictly less than $I_\eps(U)$; hence $I_\eps(U)>\Psi_\eps(\tau)$.
\end{proof}

\begin{lemma}[Interiority]\label{lem:fixed-interior}
Every minimizer $W$ of~\eqref{eq:fixed-problem} that satisfies the constraint
qualification satisfies $\gamma\le W\le1-\gamma$ for some $\gamma>0$.
\end{lemma}

\begin{proof}
Fix $\eta$, $U_1$, $U_2$ as in Definition~\ref{def:constraint-qualification}, and
write $\mathbf C(V):=(e(V),t(H,V))$.  Suppose that $W$ is not essentially
bounded away from $0$.  For large $\Lambda>0$ with $e^{-\Lambda}<\eta$, the symmetric set
$E:=\{W<e^{-\Lambda}\}$ has positive measure and is disjoint from the supports of
$U_1,U_2$.  Choose $(c_1,c_2)$ with
$\Gamma_W(U_1,U_2)(c_1,c_2)^\top=-D\mathbf C(W)[\mathbf 1_E]$.  Since
$0\le\mathcal D_W\le m$, we have $|c_1|+|c_2|\le C|E|$ with $C$ independent of
$\Lambda$, and $Z:=\mathbf 1_E+c_1U_1+c_2U_2$ satisfies $D\mathbf C(W)[Z]=0$.  The map
$(\theta,\xi_1,\xi_2)\mapsto\mathbf C(W+\theta Z+\xi_1U_1+\xi_2U_2)$ is
polynomial with invertible derivative in $(\xi_1,\xi_2)$ at the origin, so the
implicit function theorem gives $\xi_i(\theta)=O_\Lambda(\theta^2)$ with
$\mathbf C(V_\theta)=\mathbf C(W)$ for
$V_\theta:=W+\theta Z+\xi_1(\theta)U_1+\xi_2(\theta)U_2$.  For small $\theta>0$,
$V_\theta$ is a graphon.  If $0\le z<e^{-\Lambda}$ and $0<\theta<e^{-\Lambda}$, then
$J_\eps(z+\theta)-J_\eps(z)\le-\theta(\Lambda-C_\eps)$ for a constant $C_\eps$
independent of $\Lambda$, while $J_\eps'$ is bounded on $[\eta/2,1-\eta/2]$.
Therefore
\[
  I_\eps(V_\theta)-I_\eps(W)\le-\theta(\Lambda-C_\eps)|E|+C'\theta(|c_1|+|c_2|)+O_\Lambda(\theta^2)
  \le-\theta(\Lambda-C'')|E|+O_\Lambda(\theta^2),
\]
which is negative for $\Lambda$ large and then $\theta>0$ small, contradicting
minimality.  The case of values near $1$ is analogous, with $E:=\{W>1-e^{-\Lambda}\}$
and $-\mathbf 1_E$ in place of $\mathbf1_E$.
\end{proof}

\begin{lemma}[Positive multiplier]\label{lem:fixed-multiplier}
Let $W$ be a minimizer of~\eqref{eq:fixed-problem} that satisfies the
constraint qualification.  There are $\alpha\in\RR$ and $\mu>0$ such that
\begin{equation}\label{eq:fixed-el}
  J_\eps'\bigl(W(x,y)\bigr)=\alpha+\mu\,\mathcal D_W(x,y)
  \qquad\text{for almost every }(x,y).
\end{equation}
\end{lemma}

\begin{proof}
By Lemma~\ref{lem:fixed-interior}, $I_\eps$ is continuously differentiable on a
small $L^\infty$-ball around $W$ consisting of graphons, and the constraint
qualification says that $D\mathbf C(W)$ is surjective onto $\RR^2$.
Lemma~\ref{lem:lagrange} gives $\alpha,\mu\in\RR$ with
$DI_\eps(W)=\alpha\,De(W)+\mu\,Dt(H,W)$, and~\eqref{eq:fixed-el} follows as in
the proof of Lemma~\ref{lem:upper-tail-multiplier}.  Since the rows of
$\Gamma_W(U_1,U_2)$ are linearly independent, some linear combination $V$ of
$U_1,U_2$ satisfies $\int V=0$ and $\int\mathcal D_WV>0$.  Then $e(W+\theta V)=\eps$
and $t(H,W+\theta V)>\tau$ for small $\theta>0$, so
Lemma~\ref{lem:fixed-monotone} gives $DI_\eps(W)[V]\ge0$, that is,
$\mu\int\mathcal D_WV\ge0$ and $\mu\ge0$.  If $\mu=0$, then $J_\eps'(W)$ is
constant, so $W$ is constant; the edge constraint forces $W\equiv\eps$, which
contradicts $\tau>\eps^m$.
\end{proof}

\subsection{Inserting a block of vertices}

\begin{lemma}[Insertion inequality]\label{lem:fixed-insertion}
Let $W$ be a minimizer of~\eqref{eq:fixed-problem} that satisfies the
constraint qualification, and let $\alpha,\mu$ be as in
Lemma~\ref{lem:fixed-multiplier}.  Let $\mathcal F_W$ be the row
functional~\eqref{eq:row-functional} with $q=\eps$.  Then
\[
  \mathcal F_W(u)\ge c_W:=2I_\eps(W)-2\alpha\eps-\mu N\tau
  \qquad\text{for every measurable }u:[0,1]\to[0,1],
\]
and $\mathcal F_W(W_x)=c_W$ for almost every $x\in[0,1]$.
\end{lemma}

\begin{proof}
Fix $u$, and insert a block of measure $\theta$ with row $u$ as in the proof of
Lemma~\ref{lem:upper-tail-insertion}, with the value $\eps$ on the new block
and with the old values replaced by $W+\xi_1U_1+\xi_2U_2$, where $U_1,U_2$ are
as in Definition~\ref{def:constraint-qualification}.  Call the result
$W_{\theta,\xi}$.  Its edge density and $H$-density are polynomials in
$(\theta,\xi_1,\xi_2)$, with
\[
  \partial_\theta e(W_{\theta,0})\big|_{\theta=0}=2\Bigl(\int_0^1u-\eps\Bigr),
  \qquad
  \partial_\theta t(H,W_{\theta,0})\big|_{\theta=0}=\mathcal R_W(u)-N\tau,
\]
and their derivative in $(\xi_1,\xi_2)$ at the origin is
$\Gamma_W(U_1,U_2)$.  The implicit function theorem gives
$\xi(\theta)$ with $\xi(0)=0$ such that $W_\theta:=W_{\theta,\xi(\theta)}$
satisfies $e(W_\theta)=\eps$ and $t(H,W_\theta)=\tau$, and
\[
  De(W)[\xi_1'(0)U_1+\xi_2'(0)U_2]=-2\Bigl(\int_0^1u-\eps\Bigr),
  \qquad
  Dt(H,W)[\xi_1'(0)U_1+\xi_2'(0)U_2]=-\bigl(\mathcal R_W(u)-N\tau\bigr).
\]
For small $\theta>0$, $W_\theta$ is a graphon because the $U_i$ are supported
where $\eta<W<1-\eta$.  As before,
\[
  I_\eps(W_\theta)=(1-\theta)^2\int J_\eps\bigl(W+\xi_1(\theta)U_1+\xi_2(\theta)U_2\bigr)
  +2\theta(1-\theta)\int_0^1J_\eps(u)+\theta^2J_\eps(\eps),
\]
and using $DI_\eps(W)=\alpha\,De(W)+\mu\,Dt(H,W)$ we obtain
\[
  0\le\frac{\dd}{\dd\theta}\Big|_{\theta=0^+}I_\eps(W_\theta)
  =2\int_0^1J_\eps(u)-2I_\eps(W)-2\alpha\Bigl(\int_0^1u-\eps\Bigr)
  -\mu\bigl(\mathcal R_W(u)-N\tau\bigr)
  =\mathcal F_W(u)-c_W.
\]
Finally, $\int_0^1\int_0^1J_\eps(W_x(y))\dd y\dd x=I_\eps(W)$,
$\int_0^1\int_0^1W_x(y)\dd y\dd x=\eps$, and~\eqref{eq:rooted-average} give
$\int_0^1\mathcal F_W(W_x)\dd x=c_W$, so equality holds for almost every $x$.
\end{proof}

\begin{proof}[Proof of Theorem~\ref{thm:fixed-densities}]
Let $W$ be a maximizer of~\eqref{eq:entropy-maximization} that satisfies the
constraint qualification; equivalently, a minimizer of~\eqref{eq:fixed-problem}.
Apply Theorem~\ref{thm:row-criterion} with $q=\eps$, $\omega=\gamma$ from
Lemma~\ref{lem:fixed-interior}, and $\alpha,\mu$ from
Lemma~\ref{lem:fixed-multiplier}; its hypotheses are provided by
Lemma~\ref{lem:fixed-insertion}.
\end{proof}

\subsection{Degenerate maximizers occur only for a null set of values}

\begin{lemma}\label{lem:degenerate-forms}
Let $W$ be a degenerate minimizer of~\eqref{eq:fixed-problem}.  Then either
the set $\{W=0\}\cup\{W=1\}$ has positive measure, or $0<W<1$ almost everywhere
and $\mathcal D_W$ is almost everywhere equal to a constant.
\end{lemma}

\begin{proof}
Suppose that $0<W<1$ almost everywhere, and let $\eta>0$ be such that
$A:=\Omega_\eta(W)$ has positive measure.  With $\bar D:=|A|^{-1}\int_A\mathcal D_W$,
$U_1:=\mathbf 1_A$, and $U_2:=(\mathcal D_W-\bar D)\mathbf 1_A$, we have
$\det\Gamma_W(U_1,U_2)=|A|\int_A(\mathcal D_W-\bar D)^2$.  Degeneracy forces
this determinant to vanish, so $\mathcal D_W$ is constant on $\Omega_\eta(W)$.
Letting $\eta\downarrow0$, $\mathcal D_W$ is constant almost everywhere.
\end{proof}

\begin{lemma}\label{lem:degenerate-slope}
If~\eqref{eq:fixed-problem} has a degenerate minimizer at
$\tau\in(\eps^m,\tau_{\max}(\eps))$, then
\[
  \limsup_{\tau'\uparrow\tau}\frac{\Psi_\eps(\tau)-\Psi_\eps(\tau')}{\tau-\tau'}=\infty.
\]
\end{lemma}

\begin{proof}
Let $W$ be a degenerate minimizer at $\tau$ and put $W_\theta:=(1-\theta)W+\theta\eps$
for $\theta\in[0,1]$.  Then $e(W_\theta)=\eps$ and $I_\eps(W_\theta)<I_\eps(W)$
for $\theta\in(0,1)$, so the second claim of Lemma~\ref{lem:fixed-monotone} gives
$\tau_\theta:=t(H,W_\theta)<\tau$.  Since $\tau_\theta\to\tau$ as
$\theta\downarrow0$,
\[
  \frac{\Psi_\eps(\tau)-\Psi_\eps(\tau_\theta)}{\tau-\tau_\theta}
  \ge\frac{I_\eps(W)-I_\eps(W_\theta)}{\tau-\tau_\theta},
\]
and it suffices to show that the right-hand side tends to infinity along
$\theta\downarrow0$.  By convexity,
$J_\eps((1-\theta)z+\theta\eps)\le(1-\theta)J_\eps(z)$ for every $z$.

If $\{W=0\}\cup\{W=1\}$ has positive measure, then, since
$J_\eps(0)-J_\eps(\theta\eps)$ and $J_\eps(1)-J_\eps(1-\theta(1-\eps))$ are both at
least $c\,\theta\log(1/\theta)$ for small $\theta$ and some $c>0$, we get
$I_\eps(W)-I_\eps(W_\theta)\ge c'\theta\log(1/\theta)$, while
$\tau-\tau_\theta\le C\theta$ because $\theta\mapsto\tau_\theta$ is a polynomial.
Otherwise, by Lemma~\ref{lem:degenerate-forms}, $\mathcal D_W$ is a constant
$c_0$, so
$\frac{\dd}{\dd\theta}\tau_\theta\big|_{\theta=0}=\int\mathcal D_W(\eps-W)=c_0(\eps-e(W))=0$,
and hence $\tau-\tau_\theta\le C\theta^2$, while
$I_\eps(W)-I_\eps(W_\theta)\ge\theta I_\eps(W)$ with $I_\eps(W)>0$.  In both cases
the ratio tends to infinity.
\end{proof}

\begin{proof}[Proof of Corollary~\ref{cor:almost-every-tau}]
By Lemma~\ref{lem:fixed-monotone}, $\Psi_\eps$ is monotone on
$(\eps^m,\tau_{\max}(\eps))$, so by Lebesgue's theorem it is differentiable with a
finite derivative at almost every point.  At such a point, the left difference
quotients converge to the finite derivative, so by
Lemma~\ref{lem:degenerate-slope} no degenerate minimizer exists there.  Hence
the set of values $\tau$ admitting a degenerate maximizer has measure zero, and
for every other $\tau$ Theorem~\ref{thm:fixed-densities} applies to every
maximizer.
\end{proof}

\section{Conclusion}\label{sec:conclusion}

We have shown that every minimizer of a joint upper-tail problem with
Erdős–Rényi reference probability $p$ and constraints $t(H_j,W)\ge r_j^{m_j}$ for
$1\le j\le s$ is multipodal with ordered blocks. For entropy maximization
with fixed densities $e(W)=\eps$ and $t(H,W)=\tau$, the same conclusion
holds throughout the open oversaturated range
$\eps^m<\tau<\tau_{\max}(\eps)$, provided $\Delta(H)\ge2$.
Vertex insertions connect the two problems to global minimization of
$\mathcal F_W$ and $\mathcal F_{W,\eps,\tau}$, respectively; the pointwise ordering
and finiteness of their minimizers yield the ordered block structure.
For fixed densities, we introduce a constraint
qualification that establishes multipodality with ordered
blocks.  We then prove that the condition holds for optimizers
at every subgraph density in the oversaturated range.

\paragraph{Block structure and phase transitions.}  Our finiteness argument does not
control the number of blocks in the multipodal optimizers.
Previous work has established bounds on the number of blocks for specific
constrained graphs and parameter regions~\cite{ImKimLimNamYangBipodal,
KenyonRadinRenSadun2017,kenyon2014entropy,NeemanRadinSadun2023,NeemanRadinSadun2024,
RadinSadunEdge2Star,RadinSadunNearExtreme}.
A next step is to establish such bounds in more general settings.
For example, even for triangle upper tails, it remains open whether
two blocks suffice throughout the symmetry-breaking region.

For $d$-regular graphs with $d\ge2$, Lubetzky and
Zhao~\cite{lubetzky2015large} determined the boundary between constant and
nonconstant optimizers in the $(p,r)$-plane.
Determining this boundary for general $H$, and locating transitions
between different block counts would refine
our structural result. Even when the block count stays fixed, numerical
experiments for triangle upper tails suggest jumps in the block
parameters~\cite{ImKimLimNamYangBipodal}. Proving such transitions and
describing how the optimizing block parameters vary within and between
phases remain further questions.

\paragraph{Several fixed subgraph densities.}  The multipodality conclusion
for joint upper tails cannot extend to arbitrary finite families of
prescribed subgraph densities: the universality result of Cooper,
Kr\'al', and Martins~\cite{CooperKralMartins2018} allows such constraints
to force a non-multipodal entropy maximizer.
Our theorem for one fixed $H$-density in the open oversaturated range
leaves a natural question between these two settings. Which restrictions
on the number or form of the constrained graphs ensure multipodality,
and which allow non-multipodal optimizers?

\paragraph{Lower tails.}  A related question is whether every minimizer of
\[
  \min\bigl\{I_p(W):W\in\W_0,\ t(H,W)\le r^m\bigr\},
  \qquad 0<r<p<1,
\]
or its joint version with several subgraph densities, is multipodal.
For triangles, the limiting problem as $p\to0$ with $r/p$ fixed has only
constant or symmetric bipodal minimizers~\cite{HuangLi2026LowerTail}.
With fixed edge and triangle densities, symmetric bipodal phases also
occur in parts of the undersaturated region $\tau<\eps^3$ in the
$(\eps,\tau)$-plane~\cite{NeemanRadinSadun2024}.
In these phases, connections are stronger between the blocks than within
them, so the two row types cross. Establishing multipodality for general
lower-tail constraints at fixed $p$, or throughout the fixed-density
undersaturated region, would therefore require an argument that allows
crossing rows.


\paragraph{AI disclosure.}
\red{[To be written by the authors: describe the role of AI assistance in the
results, proofs, and writing of this paper, in the same format as the
companion paper.]}

\paragraph{Acknowledgments.}
% TODO: confirm that the grants below apply to this paper.
S. Im was supported by a KIAS individual Grant (AP109501) at Korea Institute for Advanced Study.
S. Lim, T. Kim, and H. Yang were supported by the National Research Foundation of Korea (NRF)
grant funded by the Korean Government (MSIT) (No.~RS-2023-00279680).

\bibliographystyle{plain}
\bibliography{cites}

\appendix

\section{Comparison with Alvarado, Ramon, and Wang}\label{app:arw}

Alvarado, Ramon, and Wang (ARW)~\cite{AlvaradoRamonWang2026} develop a
differential geometric approach to entropy optimization under general
multi-relational equality constraints. Our results treat joint upper
tails and entropy optimization with a fixed edge density and one fixed
subgraph density in the open oversaturated range. We compare ARW's
assumptions with those of \Cref{thm:joint-upper-tail,thm:fixed-densities}.
Then we present a construction suggesting that additional assumptions may be needed
to justify the finite-step conclusion for all global minimizers
in~\cite[Theorem~5.25]{AlvaradoRamonWang2026}.

\subsection{Comparison of assumptions}\label{app:arw-setting}

We state ARW's conditions for a single constraint $C(W)=u$, where
\[
  C:L^\infty_{\mathrm{sym}}([0,1]^2)\longrightarrow\RR,
  \qquad C(W):=\sum_{i=1}^q c_i t(F_i,W).
\]
Here $q\ge1$, $F_1,\dots,F_q$ are finite simple graphs, $c_i\in\RR$, and
$u\in\RR$ is the prescribed constraint value.
For an integer $k\ge1$, a \emph{$k$-step representation} consists
of a symmetric real matrix $A=(a_{ij})_{i,j=1}^k$ and a mass vector
$\pi=(\pi_1,\dots,\pi_k)$ with $\pi_i\ge0$ and $\sum_{i=1}^k\pi_i=1$.
Partitioning $[0,1]$ into consecutive intervals $S_i$ of length $\pi_i$ gives
\[
  W_{A,\pi}:=\sum_{i,j=1}^k a_{ij}\1_{S_i\times S_j}.
\]
The map $(A,\pi)\mapsto C(W_{A,\pi})$ is polynomial, and a representation
is \emph{feasible} when $C(W_{A,\pi})=u$.

With $\pi$ fixed, use the $k(k+1)/2$ independent entries $a_{ij}$, $i\le j$,
as coordinates on the space of symmetric matrices. The Jacobian is the row vector
\[
  D_A\bigl[C(W_{A,\pi})\bigr]:=
  \left(\frac{\partial C(W_{A,\pi})}{\partial a_{ij}}\right)_{1\le i\le j\le k}
  \in\RR^{1\times k(k+1)/2}.
\]
The value $u$ is \emph{regular}~\cite[Definition~5.5]{AlvaradoRamonWang2026}
if this row vector is nonzero at every feasible real-valued representation,
for every $k$.

For non-extremality, let $T_k$ be the set of constraint values attained
by $k$-step graphons:
\[
  T_k:=\left\{C(W_{A,\pi}):
    A=A^\top,\ 0\le a_{ij}\le1,\ \pi_i\ge0,\
    \sum_{i=1}^k\pi_i=1\right\}.
\]
A regular value $u$ is \emph{non-extremal} if
$u\in T_k\setminus\bigcup_{\ell\ge1}\partial T_\ell$ for some $k$.
This excludes boundaries of the attainable density sets at these block
sizes.

For a real-valued objective $I$, ARW's \emph{independence}
hypothesis~\cite[Definition~5.4 and Theorem~5.25]{AlvaradoRamonWang2026}
requires that for every $k\ge2$ there is a representation with $\pi_i>0$
for all $i$ at which both parameterized functions are analytic and
\[
  \operatorname{rank}
  \begin{pmatrix}
    D_A\bigl[C(W_{A,\pi})\bigr]\\
    D_A\bigl[I(W_{A,\pi})\bigr]
  \end{pmatrix}=2.
\]

Finally, a global minimizer $W$ of $I$ under $C(W)=u$ is \emph{isolated}
in cut distance if there is $\rho>0$ such that every feasible graphon $V$ satisfies
\[
  0<\delta_\cut(W,V)<\rho\quad\Longrightarrow\quad I(V)>I(W).
\]
Equivalently, no other global minimizing class lies within cut distance
$\rho$ of $W$.
Under the assumptions above, including isolation of every global minimizing
class, the result in~\cite[Theorem~5.25]{AlvaradoRamonWang2026} asserts that
all global entropy minimizers admit finite-step representations.

For the two constraint problems studied in this paper,
\Cref{thm:joint-upper-tail,thm:fixed-densities} do not assume ARW's analytic
independence, non-extremal regularity, or isolation.
For fixed edge and subgraph densities, we use local independence of their
derivatives at an optimizer $W$, expressed by the constraint
qualification in \Cref{def:constraint-qualification}:
\[
  \det\Gamma_W(U_1,U_2)\ne0
\]
for bounded symmetric $U_1,U_2$ supported in $\Omega_\gamma(W)$ for some
$\gamma>0$. We prove this for every optimizer in the open oversaturated
range (\Cref{prop:optimizer-nondegeneracy}), without asserting regularity
at all feasible representations. For joint upper tails,
\Cref{lem:upper-tail-multiplier} obtains multipliers by separation of convex
sets. We do not assume isolation of optimizing graphons;
our results do not exclude a sequence of distinct optimizing classes
converging to $W$ in cut distance.

\subsection{A construction testing the finite-step conclusion}\label{app:arw-entropy-counterexample}

Let
$W_{\mathrm h}(x,y):=\1_{\{x+y\le1\}}$ be the \emph{half-graphon}, and set
\[
  W_*:=\frac14+\frac12W_{\mathrm h}.
\]
Lov\'asz and Szegedy~\cite{LovaszSzegedy2011} show that the half-graphon is
finitely forcible among bounded real-valued symmetric kernels and that
affine transformations preserve this property. Thus $W_*$ is
\emph{finitely forcible} in this class:
there are finite simple graphs $F_1,\dots,F_\ell$ such that, for every
bounded symmetric kernel $U$,
\[
  t(F_i,U)=t(F_i,W_*)\quad(1\le i\le\ell)
  \quad\Longleftrightarrow\quad U\sim W_*.
\]
Here \emph{weak isomorphism}, denoted by $U\sim W$, means equality of
$t(F,U)$ and $t(F,W)$ for every finite simple graph $F$.
For graphons it is equivalent to having zero cut distance; write $[W]$ for
the equivalence class.

Define
\[
  q(W):=\sum_{i=1}^\ell\bigl(t(F_i,W)-t(F_i,W_*)\bigr)^2,
  \qquad
  g(W):=\bigl(e(W)-\tfrac14\bigr)q(W).
\]
Since
\[
  t(F\sqcup F',W)=t(F,W)t(F',W),\qquad t(K_1,W)=1,
\]
where $\sqcup$ denotes disjoint union and $K_1$ is the one-vertex graph,
$g$ is a finite real linear combination of graph densities.
We take $C=g$ and $u=0$, so the constraint is $g(W)=0$.
Up to a positive constant factor, ARW's entropy objective is
\[
  I(W):=\int_{[0,1]^2}I_0(W(x,y))\dd x\dd y,
  \qquad I_0(z):=z\log z+(1-z)\log(1-z),
\]
where $0\le z\le1$ and $0\log0:=0$.

\begin{prop}\label{prop:arw-entropy-counterexample}
The constraint $g(W)=0$ satisfies the independence and non-extremal
regularity hypotheses of~\cite[Theorem~5.25]{AlvaradoRamonWang2026},
including independence with the entropy objective.  Its global
entropy minimizers are exactly the two isolated classes $[1/4]$ and
$[W_*]$, and $W_*$ is not multipodal.
\end{prop}

\begin{proof}
The feasible graphons $W$ satisfy either $e(W)=1/4$ or $W\sim W_*$ by construction.
If $e(W)=1/4$, Jensen's inequality and the strict convexity of $I_0$ give
\[
  I(W)\ge I_0\bigl(e(W)\bigr)=I_0(1/4),
\]
with equality if and only if $W=1/4$ almost everywhere.
Weakly isomorphic graphons $U,W$ admit measure-preserving maps $\phi,\psi$
such that $U^\phi=W^\psi$ almost everywhere, so $I(U)=I(W)$.
Consequently, if $W\sim W_*$, then
\[
  I(W)=I(W_*)=\tfrac12 I_0(1/4)+\tfrac12 I_0(3/4)=I_0(1/4).
\]
Thus there are exactly two minimizing classes.
Since $W_*-\tfrac14=\tfrac12W_{\mathrm h}\ge0$, its cut norm equals its
integral. Relabeling preserves this integral and leaves the constant
graphon unchanged, so
\[
  \delta_\cut\bigl(\tfrac14,W_*\bigr)
  =\bigl\|W_*-\tfrac14\bigr\|_\cut
  =\frac12\int_{[0,1]^2}W_{\mathrm h}(x,y)\dd x\dd y=\frac14.
\]
Both minimizing classes are therefore isolated.
The degree distribution, the law of $d_W(X)$ for uniform $X\in[0,1]$,
is invariant under weak isomorphism because its moments are star densities.
For $W_*$,
\[
  d_{W_*}(x)=\frac34-\frac{x}{2},
\]
so this distribution is uniform on $[1/4,3/4]$.
A finite-step kernel has only finitely many degree values and
therefore cannot be weakly isomorphic to $W_*$.

Every feasible real-valued finite-step kernel $U$ consequently has
$q(U)>0$ and $e(U)=1/4$.  Writing
$Dg(U)[V]:=\left.\frac{\dd}{\dd t}g(U+tV)\right|_{t=0}$ for bounded
symmetric $V$, we obtain
\[
  Dg(U)[V]=q(U)\int_{[0,1]^2}V,
  \qquad Dg(U)[1]=q(U)>0.
\]
The constant direction is available at every step size, so $0$ is
regular.  Along constant graphons $c$ near $1/4$,
\[
  g(1/4)=0,\qquad
  \left.\frac{\dd}{\dd c}g(c)\right|_{c=1/4}=q(1/4)>0.
\]
Hence $0\in\operatorname{int}T_k$ for every $k$, with entries
restricted to $[0,1]$, and $0$ is non-extremal.

It remains to verify independence of $g$ and $I$: for every $k\ge2$,
we need a $k$-step representation at which the Jacobian of $(g,I)$ in
the entries of $A$ has rank $2$.  For $k=2$, take $\pi=(1/2,1/2)$ and
\[
  A_\delta:=\begin{pmatrix}1/4+\delta&1/4\\1/4&1/4-\delta\end{pmatrix},
  \qquad U_\delta:=W_{A_\delta,\pi},\qquad 0<\delta<1/4.
\]
With $\pi$ fixed, the Jacobian has two rows, for $g$ and $I$, and three
columns, for $a_{11},a_{12},a_{22}$.  We choose the columns
for $a_{11}$ and $a_{22}$.
Since $e(U_\delta)=1/4$, these columns form the matrix
\[
  \frac{\partial(g,I)}{\partial(a_{11},a_{22})}(A_\delta,\pi)
  =\frac14
  \begin{pmatrix}
    q(U_\delta)&q(U_\delta)\\
    I_0'(1/4+\delta)&I_0'(1/4-\delta)
  \end{pmatrix}.
\]
Its determinant is
\[
  \det\frac{\partial(g,I)}{\partial(a_{11},a_{22})}(A_\delta,\pi)
  =\frac{q(U_\delta)}{16}
     \bigl(I_0'(1/4-\delta)-I_0'(1/4+\delta)\bigr)\ne0,
\]
since $q(U_\delta)>0$ and $I_0'$ is strictly increasing.
Thus the full Jacobian has rank $2$.
Both functions are analytic at these interior entries.

For $k>2$, refine the two blocks into $k$ positive-mass subblocks.
Consider the two parameter directions that vary all refined entries
within the first or the second original diagonal block together.
The derivative of $(g,I)$ in these directions is the same nonsingular
$2\times2$ matrix as above, so the full refined Jacobian also has rank $2$.
All entries remain in $(0,1)$, preserving analyticity and proving
independence for every $k\ge2$.
\end{proof}

Thus the assumptions stated in~\cite[Theorem~5.25]{AlvaradoRamonWang2026}
do not exclude an isolated global minimizer with no finite-step representation.

This example suggests extending the regularity requirement to all global
optimizers, including those without finite-step representations.
Each density difference defining $q$ vanishes at $W_*$, so for every
bounded symmetric $V$,
\[
  q(W_*+tV)=O(t^2),\qquad g(W_*+tV)=O(t^2),\qquad Dg(W_*)[V]=0.
\]
Requiring a nonzero constraint derivative at every global optimizer would
therefore exclude this example; for several constraints, the corresponding
condition is surjectivity of their joint derivative.
A separate argument would be needed to establish that this stronger
condition suffices for the general finite-step conclusion.


\end{document}
